% The Rust chapter of the WFDB-x user's manual. It follows the sections of the Python
% chapter (chapters/python.tex), in the same order. Each \apientry of the section "The
% high-level API" is a public free function of the crate root defined in
% rust/src/record.rs, rust/src/check.rs or rust/src/metadata.rs (each also a method of
% Wfdb, with the same parameters after &mut self); its signature listing is copied from
% that definition, and its params name every parameter of the definition, in order. The
% examples of style "example" were compiled and run (each as the body of main(), after
% the lines given in the conventions of the section, in a copy of tests/data);
% check_manual.py does not run them.
\chapter{Rust}
\label{chap:rust}
\chapterlanguage{rust}

The Rust implementation of WFDB-x is the crate \code{wfdbx}, in the directory \file{rust}
of the repository: a library and the 33 command-line tools, each a separate program of
the crate. It is a complete Rust implementation of the WFDB Software Package: it uses only
the Rust standard library (no other crate), contains no \code{unsafe} code, and never
calls the C library.

It offers WFDB-x in three ways:
\begin{itemize}
\item \textbf{The high-level API} (\cref{sec:rust-api}): functions that read and write
  whole records, headers, annotation files and header metadata as \code{Vec}s and
  structures, in one call, with options given as fields of a structure, and that return a
  \code{Result} whose error type is \code{wfdbx::record::Error}. They read signal files in
  large blocks, and are the fastest way to read a record. They are an addition of WFDB-x:
  the C package has no equivalent. They are the recommended interface for new code and
  analysis.
\item \textbf{The WFDB library functions} under their C names (\code{isigopen},
  \code{getvec}, \code{getann}, \code{putann}, \ldots), with the behavior of the C library
  (\cref{sec:rust-library}). They are documented once, for all the languages, in
  \cref{chap:library}.
\item \textbf{The command-line tools} (\code{rdsamp}, \code{rdann}, \code{wrsamp},
  \code{xform} and 29 others), with the options, output and exit status of the C tools,
  one program per tool, run by its name (\cref{sec:rust-tools}). They are documented in
  \cref{chap:tools}.
\end{itemize}
The high-level functions return exactly the samples and annotations that the library
functions (and the C library) give.

% ==========================================================================================
\section{Installation and requirements}
\label{sec:rust-install}

\subsection{Requirements}

WFDB-x for Rust needs Rust 1.70 or later (the stable toolchain, edition 2021; the crate
declares \code{rust-version = "1.70"}) and Cargo, on Linux, macOS or Windows. It has no
dependency beyond the Rust standard library: no WFDB library and no libFLAC (the FLAC
formats 508, 516 and 524 are read and written by WFDB-x's own codec). Network access
(\cref{sec:network}) runs the program \cmd{curl}, which must then be installed. The tests
of \code{check\_record} also run \cmd{python3}, with NumPy (\cref{sec:rust-testing}).

\subsection{Using the library in a Cargo project}

A Cargo project uses the crate as a dependency, from the repository or from a copy of it,
in its \file{Cargo.toml}:

\begin{lstlisting}[language=TOML]
[dependencies]
wfdbx = { git = "https://github.com/alphanumericslab/wfdb-x" }
# or, from a copy of the repository:
wfdbx = { path = "../wfdb-x/rust" }
\end{lstlisting}

The crate is \code{wfdbx}, version \code{1.2.1}: its library under the GNU Lesser
General Public License, version 2.1 or later, and its tools (the binaries of
\file{src/bin}) under the GNU General Public License, version 3 or later
(\cref{sec:license}). Programs that use it are built in release mode for speed
(\cmd{cargo build -{}-release}, \cmd{cargo run -{}-release}).

\subsection{Building and installing the tools}

From the root of a copy of the repository, \cmd{cargo install} builds the tools in
release mode and installs them in \file{\textasciitilde/.cargo/bin}; \cmd{cargo build
-{}-release}, in the directory \file{rust}, builds the library and the tools in
\file{rust/target/release/} without installing them:

\begin{lstlisting}[style=shell]
cargo install --path rust                                # the 33 tools, in ~/.cargo/bin
cargo build --release --manifest-path rust/Cargo.toml    # in rust/target/release/
cargo build --release --manifest-path rust/Cargo.toml --bin rdsamp   # one tool
\end{lstlisting}

The tools, like the C tools, have no \opt{-{}-version} option. \cmd{cargo doc -{}-open}, in
the directory \file{rust}, builds the documentation of the crate (its doc comments, with
examples) and opens it.

\subsection{The sample record}
\label{sec:rust-resources}

The repository includes the sample record \code{100s}, the sample record of the WFDB
Software Package: the first minute of record 100 of the MIT-BIH Arrhythmia Database, two
ECG signals (\code{MLII} and \code{V5}) sampled at 360\,Hz, with the reference beat
annotations (annotator \code{atr}). Its files (\file{100s.hea}, \file{100s.dat},
\file{100s.atr}) are in the directory \file{tests/data} of a copy of the repository, with
the other test records of WFDB-x. The crate does not install them: programs name that
directory as their input directory (\code{wfdbx\_set\_input\_dir("tests/data")}, or the
option \code{input\_dir}). Work that uses the MIT-BIH Arrhythmia Database cites it as
described on PhysioNet.

% ==========================================================================================
\section{Quick start}
\label{sec:rust-quickstart}

This program reads the header, ten seconds of the signals and the annotations of the
sample record \code{100s}, from the directory \file{tests/data} of the repository (the
input directory of each call), and computes the mean heart rate. It then writes the first
ten seconds as a new record, \code{first10}, in the current directory, with a metadata
entry in its header.

\begin{lstlisting}[language=Rust,style=example]
use wfdbx::{check_record, isqrs, read_annotations_with, read_header_with, read_metadata,
            read_record, set_metadata, write_record, ReadOptions, WriteOptions};

fn main() -> Result<(), Box<dyn std::error::Error>> {
    // The sample record 100s, in the directory tests/data of the repository.
    let sample = ReadOptions { input_dir: Some("tests/data".into()), ..Default::default() };

    // The header only: no samples are read.
    let hdr = read_header_with("100s", &sample)?;
    println!("{} {} {}", hdr.nsig, hdr.fs, hdr.nsamples); // 2 360 21600

    // Seconds 10 to 20 of both signals, frame by frame, in ADC units.
    let opt = ReadOptions {
        from: Some("0:10".into()),
        to: Some("0:20".into()),
        ..sample.clone()
    };
    let rec = read_record("100s", &opt)?;
    println!("{} {}", rec.nsamples(), rec.nsig); // 3600 2
    println!("{} {}", rec.at(0, 0), rec.physical(0, 0)?); // 946 -0.39

    // The beat annotations, and the mean heart rate.
    let ann = read_annotations_with("100s", "atr", None, &sample)?;
    let beats: Vec<i64> = (0..ann.len())
        .filter(|&i| isqrs(ann.codes[i]) != 0)
        .map(|i| ann.samples[i])
        .collect();
    let rr = (beats[beats.len() - 1] - beats[0]) as f64 / ann.fs / (beats.len() - 1) as f64;
    println!("{} beats, {:.1} beats per minute", beats.len(), 60.0 / rr);
    // 74 beats, 73.9 beats per minute

    // Write the first 10 seconds as the record "first10" (first10.hea and first10.dat,
    // in the current directory), with the signal specifications of 100s.
    let opt = ReadOptions { to: Some("0:10".into()), ..sample.clone() };
    let x = read_record("100s", &opt)?;
    let w = WriteOptions {
        gains: vec![200.0],
        baselines: vec![1024],
        adczeros: vec![1024],
        adcres: vec![11],
        units: vec![Some("mV".into())],
        names: x.names(),
        metadata: vec![("Source".into(), "MIT-BIH record 100".into())],
        ..Default::default()
    };
    write_record("first10", &x.samples, x.nsig, x.fs, &w)?;

    // Header metadata: "# Key: value" lines of the header (keys and values are bytes).
    set_metadata("first10", "Age", "69")?;
    for (key, value) in read_metadata("first10")? {
        println!("{} = {}", String::from_utf8_lossy(&key), String::from_utf8_lossy(&value));
    }
    // Source = MIT-BIH record 100
    // Age = 69
    println!("{}", check_record("first10", &[])?); // no problems found
    Ok(())
}
\end{lstlisting}

Run from the root of the repository, as the program \file{src/main.rs} of a Cargo project
that depends on \code{wfdbx} (\cmd{cargo run -{}-release}). The header \file{100s.hea}
gives no units: the program writes \code{"mV"} explicitly, since WFDB-x never assumes
units, a gain, an ADC resolution or a sampling frequency that a record does not give
(\cref{chap:behavior}).

% ==========================================================================================
\section{The high-level API}
\label{sec:rust-api}

The high-level API is made of the functions of this section, defined in the modules
\code{wfdbx::record}, \code{wfdbx::check} and \code{wfdbx::metadata} and available at the
root of the crate (\code{wfdbx::read\_record}, \ldots), of the structures of their options
(\cref{sec:rust-options}) and of their results (\cref{sec:rust-results}), and of the error
type \code{wfdbx::record::Error} (\cref{sec:rust-warnings}). The doc comments of the
crate (\cmd{cargo doc}) document each of them.

\subsection{Conventions}
\label{sec:rust-conventions}

\paragraph{Free functions and methods.} The state of the library (the open record, the
annotators, the input and output directories, network access, legacy mode) is held by a
value of the type \code{Wfdb}. Each high-level function exists twice: as a method of
\code{Wfdb}, such as \code{w.read\_record("100s", \&opt)}, which uses the state of
\code{w}, and as a free function of the crate with the same name and the same
parameters, such as \code{wfdbx::read\_record("100s", \&opt)}, which uses the default
\code{Wfdb} value of the calling thread (one per thread). The entries of this section
document the free functions; the methods take \code{\&mut self} first, and then the same
parameters, with the same types and results. Separate \code{Wfdb} values share no state
(except what belongs to the process: the count of write errors, the cache of fetched URLs
and the messages printed once per process), so several
records can be read at once, in one thread or in several; a \code{Wfdb} is \code{Send},
and is used by one thread at a time.

\begin{lstlisting}[language=Rust,style=example]
use wfdbx::Wfdb;

fn main() -> Result<(), wfdbx::record::Error> {
    let mut w = Wfdb::new(); // a library state of its own
    w.wfdbx_set_input_dir("tests/data");
    let a = w.read_record("100s", &Default::default())?; // on w
    let b = wfdbx::read_record("tests/data/100s", &Default::default())?; // on the default value
    println!("{}", a.samples == b.samples); // true
    Ok(())
}
\end{lstlisting}

\paragraph{Options.} The options of a function are the fields of a structure
(\code{ReadOptions}, \code{WriteOptions}, \code{AnnWriteOptions}, \code{MetadataOptions},
\cref{sec:rust-options}), passed by reference as the last argument. Each implements
\code{Default}, so that the fields that are not given keep their defaults:
\code{ReadOptions \{ highres: true, ..Default::default() \}}. Several functions have a
short form with no options or with a few arguments (\code{read\_header},
\code{read\_annotations}, \code{write\_annotations}, \code{check\_record}, the metadata
functions), and a form whose name ends in \code{\_with}, which takes the options
structure. An option that may be absent is an \code{Option}; \code{None} means the
default or the current setting.

\paragraph{Record names and directories.} A record is named as in the other parts of
WFDB-x (\cref{chap:records,sec:record-paths}), without the \file{.hea} extension (the
readers, and \code{write\_annotations}, also accept a trailing \code{".hea"};
\code{write\_record} does not). A record name is one of:
\begin{itemize}
\item a full path, absolute (\code{"/data/mitdb/100"}) or in the home directory
  (\code{"\textasciitilde/Documents/100s"}), used as given;
\item a relative path (\code{"100s"}, \code{"data/100s"}), read from the input directory
  if one is set, otherwise from the current directory, and written to the output
  directory if one is set, otherwise to the current directory (missing directories are
  created);
\item a URL (\code{"https://physionet.org/files/mitdb/1.0.0/100"}), read only if network
  access is on; a URL cannot be written.
\end{itemize}
The input and output directories are set with \code{wfdbx\_set\_input\_dir} and
\code{wfdbx\_set\_output\_dir} (\cref{sec:rust-settings}) or with the environment
variables \code{WFDBX\_INPUT\_DIR} and \code{WFDBX\_OUTPUT\_DIR}. The readers take the
optional fields \code{input\_dir} and \code{network} of a \code{ReadOptions} (or a
\code{MetadataOptions}), and the writers a \code{dir}: they apply to that one call, and
the previous settings are restored when it returns, also when it fails.
\code{input\_dir} applies to relative names only, not to full paths or URLs. A writer
given a name with a path creates the files in that directory, and the header names the
record and its signal file without the path (\code{rec}, \code{rec.dat}), so that the
record opens under any path that leads to it; in legacy mode, as in the original, the
header gives the name with its path. With \code{dir}, the name must be a plain name
(\code{"rec"}): the files are created in \code{dir}, and the record opens from that
directory under its plain name. The signal files that a header names are read from the
header's directory.

\paragraph{Times.} The fields \code{from} and \code{to} of \code{ReadOptions}, and the
arguments \code{from} and \code{to} of \code{read\_annotations}, are
\code{Option<TimeSpec>} values (\cref{tab:rust-timespec}): a sample number or a WFDB time
string, as \code{strtim} reads it (\code{"0:10"}, 10 seconds; \code{"1:02:03.5"};
\code{"s3600"}, sample 3600; \code{"e"}, the end of the record; or a time of day in
brackets, \code{"[14:00:00]"}). A \code{TimeSpec} converts from integers and strings, so
\code{Some(3600.into())} and \code{Some("0:10".into())} both work; \code{None} means that
no time is given. A time in seconds needs the sampling frequency; when the record does
not give it, such a time is an error.

\paragraph{Samples.} Samples are \code{i32} values (C's \code{WFDB\_Sample}), in ADC
units, stored frame by frame in one \code{Vec<i32>}: \code{samples[i * nsig + s]} is
sample \code{i} of signal \code{s}. Invalid samples are \code{WFDB\_INVALID\_SAMPLE}
($-32768$), and NaN in physical units. Sample numbers and times are \code{i64}.

\paragraph{Per-signal and per-annotation values.} The per-signal fields of
\code{WriteOptions} (\code{gains}, \code{units}, \ldots) hold one value for all the
signals, or one value per signal; the per-annotation fields of \code{Annotations}
(\code{symbols}, \code{chans}, \ldots) are empty or hold one value per annotation. Any
other length is an error, reported before anything is written.

\paragraph{Signal parameters.} No sampling frequency, gain, ADC resolution or units is
ever assumed. Where one is needed and the record does not give it, it must be given;
otherwise the function returns an error. A record whose header gives no sampling
frequency has \code{fs} 0, a signal without an ADC resolution has \code{adcres} 0, and a
signal without units has \code{units} \code{None}.

\paragraph{Text.} Text read from files (signal descriptions, units, info strings,
metadata, aux strings, mnemonics) is given as bytes, \code{Vec<u8>}, since WFDB files are
not UTF-8: every byte sequence round-trips unchanged. Arguments accept \code{\&str},
\code{\&[u8]}, \code{b"..."} or \code{Vec<u8>} (\code{impl AsRef<[u8]>}), and a
\code{Vec<u8>} field is set from a \code{\&str} with \code{.into()}.
\code{String::from\_utf8\_lossy} converts bytes for printing. The fields
\code{channels} and \code{input\_dir} of the options and \code{base\_time} of the results
are \code{String}s, and the directories \code{dir} are \code{PathBuf}s.

\paragraph{Legacy mode.} With \code{legacy} set to \code{true}, a function behaves as in
the original WFDB package, with its assumed signal parameters and its search for files
(\cref{chap:behavior}); the library's legacy mode is on for the duration of the call
only. With \code{false} (the default), the current setting of \code{wfdbx\_set\_legacy}
applies, which is off unless it was turned on.

\paragraph{State.} The functions that read and write records and annotation files call
\code{wfdbquit()} on the \code{Wfdb} value they use, before and after their work, also
when they fail, which closes anything opened with the WFDB library functions on that
value. The library's settings (legacy mode, getvec mode, input and output directories,
network access) are the same after a call as before it. \code{read\_record},
\code{read\_header}, \code{read\_annotations}, \code{write\_record},
\code{write\_record\_physical}, \code{write\_annotations}, their \code{\_with} forms and
\code{check\_record} also leave the time resolution of output annotators
(\code{getafreq()}) and the getvec mode (\code{getgvmode()}) as they were, also after
reading an EDF record, whose opening sets both (as in C; the library functions keep that
behavior). The metadata functions do not
call \code{wfdbquit()}: they open only the header, and leave the library's open files as
they are. Results are returned by value.

\paragraph{The examples} of this section are the body of the function
\code{fn main() -> Result<(), Box<dyn std::error::Error>{}>}, followed by \code{Ok(())},
compiled after these lines, and run with the free functions in a copy of the test
records of \file{tests/data} (among them \code{100s}), which is the current directory and
holds the records that the examples write:
\begin{lstlisting}[language=Rust]
use wfdbx::*;

/// Text (bytes) as a String, for printing.
fn s(b: &[u8]) -> String {
    String::from_utf8_lossy(b).into_owned()
}
\end{lstlisting}

% ------------------------------------------------------------------------------------------
\subsection{Reading records}
\label{sec:rust-reading}

\apientry{read\_record}{Read the signals of a record}
\begin{lstlisting}[language=Rust,style=signature]
pub fn read_record(record: impl AsRef<[u8]>, opt: &ReadOptions) -> Result<Record, Error>
\end{lstlisting}
Reads the signals of a record from \code{opt.from} up to \code{opt.to} (exclusive), in
ADC units, frame by frame, together with the description of the record. The samples are
those that \code{getvec} returns (\cref{chap:library}), read in large blocks with
\code{getvecs}. If the record ends before \code{opt.to}, the samples up to its end are
returned. When the read reaches the end of the record, the samples are compared with the
checksums of the header (\cref{sec:rust-warnings}). The getvec mode (\code{getgvmode()})
and the time resolution of output annotators (\code{getafreq()}) are the same after the
call as before it.
\begin{params}
\item[record] Record name, without the \file{.hea} extension (a trailing \code{".hea"} is
  also accepted): a full path, absolute (\code{"/data/mitdb/100"}) or in the home
  directory (\code{"\textasciitilde/Documents/100s"}), used as given; a relative path
  (\code{"100s"}, \code{"data/100s"}), read from the input directory if one is set,
  otherwise from the current directory; or a URL
  (\code{"https://physionet.org/files/mitdb/1.0.0/100"}), read only if network access is
  on.
\item[opt] The options (\cref{tab:rust-readoptions}): the signals to read
  (\code{channels}), the interval (\code{from}, \code{to}), \code{highres}, \code{pad},
  \code{checksum\_warning}, \code{legacy}, and the input directory (under which relative
  record names are read; it does not apply to full paths or URLs) and network access of
  the call (\code{input\_dir}, \code{network}). \code{\&ReadOptions::default()}: every
  signal of the whole record, with the current settings.
\end{params}
\returns{A \code{Record} (\cref{tab:rust-record}), with the samples in ADC units.}
\raises{An \code{Error} if the record cannot be read or has no signals (``can't read record
REC''), or its signals cannot be opened (``can't open signals of record REC''); if a
requested signal does not exist (``record REC has no signal C''); if \code{opt.from} is a
negative sample number (``read\_record: from must not be negative'') or cannot be reached
(``can't seek to sample N of record REC''); if a time string needs the sampling frequency
and the header gives none (``sampling frequency of record REC is not specified''); or if a
signal file is shorter than its header says, or cannot be read or decoded (``can't read
signal data of record REC (getvec status S)''). Data that do not match the checksums of
the header are returned, with a warning (see \code{checksum\_warning}).}
\begin{lstlisting}[language=Rust,style=example]
// Seconds 10 to 20 of lead MLII, in physical units.
let opt = ReadOptions {
    channels: vec!["MLII".into()],
    from: Some("0:10".into()),
    to: Some("0:20".into()),
    ..Default::default()
};
let rec = read_record("100s", &opt)?;
let mlii: Vec<f64> = rec.physical_signal(0)?;
println!("{} {} {} {}", mlii.len(), rec.fs, rec.sampfrom, s(&rec.names()[0]));
// 3600 360 3600 MLII

// The first 3 samples of both signals, in ADC units, V5 first.
let first = ReadOptions {
    channels: vec!["1".into(), "0".into()],
    to: Some(3.into()),
    ..Default::default()
};
let rec = read_record("100s", &first)?;
for i in 0..rec.nsamples() {
    println!("{} {}", rec.at(i, 0), rec.at(i, 1)); // 1011 995, 3 times
}
println!("{} {}", rec.siginfo[0].gain, rec.siginfo[0].baseline); // 200 1024

// The whole record: the checksums are verified.
let rec = read_record("100s", &ReadOptions::default())?;
println!("{} {}", rec.nsamples(), rec.checksum_ok); // 21600 true
\end{lstlisting}
A multi-frequency record (here the test record \code{mfreq}, whose signals have 1, 2
and 4 samples per frame at a frame frequency of 100\,Hz) is read at its frame frequency,
or at its highest sampling frequency:
\begin{lstlisting}[language=Rust,style=example]
let rec = read_record("mfreq", &ReadOptions::default())?;
println!("{} {}", rec.fs, rec.nsamples()); // 100 600
let opt = ReadOptions { highres: true, ..Default::default() };
let rec = read_record("mfreq", &opt)?;
println!("{} {}", rec.fs, rec.nsamples()); // 400 2400
let spf: Vec<i32> = rec.siginfo.iter().map(|si| si.spf).collect();
println!("{spf:?}"); // [1, 2, 4]
\end{lstlisting}

\apientry{read\_header}{Read the header of a record}
\begin{lstlisting}[language=Rust,style=signature]
pub fn read_header(record: impl AsRef<[u8]>, legacy: bool) -> Result<Header, Error>
\end{lstlisting}
Reads the header of a record, without opening its signal files. The result describes the
record and each of its signals, and holds the info strings and the metadata entries of the
header (from the \file{.hea} file only, as \code{read\_metadata} reads them, with its
warning for a key given more than once). For an EDF record, which has no \file{.hea}
file, and for a header read from the standard input (\code{"-"}), the metadata is empty,
without an error. For a multi-segment record, the master header is read, and
\code{nsamples} is the length of the whole record.
\begin{params}
\item[record] Record name, as for \code{read\_record}: a full path (absolute or
  \code{\textasciitilde}), used as given; a relative path, read from the input directory
  if one is set, otherwise from the current directory; or a URL, read only if network
  access is on.
\item[legacy] \code{true} to read the header as the original WFDB package does; normally
  \code{false}.
\end{params}
\returns{A \code{Header} (\cref{tab:rust-header}).}
\raises{An \code{Error} if the header cannot be read (``can't read header for record
REC'').}
\begin{lstlisting}[language=Rust,style=example]
let hdr = read_header("100s", false)?;
println!("{} {} {} {}", s(&hdr.name), hdr.fs, hdr.nsig, hdr.nsamples); // 100s 360 2 21600
for si in &hdr.siginfo {
    let desc = si.desc.as_deref().unwrap_or(b"");
    let fname = si.fname.as_deref().unwrap_or(b"");
    println!("{} {} {} {} {} {:?}", s(desc), s(fname), si.fmt, si.gain, si.adcres, si.units);
}
// MLII 100s.dat 212 200 11 None
// V5 100s.dat 212 200 11 None
println!("[{}]", s(&hdr.comments[0])); // [ 69 M 1085 1629 x1]
println!("{}", hdr.base_time.is_none()); // true: the header gives no start time
let (key, value) = &hdr.metadata[0];
println!("[{}] [{}]", s(key), s(value));
// [Produced by xform from record 100, beginning at 0] [0]
\end{lstlisting}

\apientry{read\_header\_with}{Read the header of a record, with options}
\begin{lstlisting}[language=Rust,style=signature]
pub fn read_header_with(record: impl AsRef<[u8]>, opt: &ReadOptions) -> Result<Header, Error>
\end{lstlisting}
As \code{read\_header}, with the options \code{input\_dir}, \code{network} and
\code{legacy} of \code{opt}, for this call only. The other fields of \code{opt} are not
used.
\begin{params}
\item[record] Record name, as for \code{read\_header}.
\item[opt] The options (\cref{tab:rust-readoptions}); only \code{input\_dir} (optional:
  the input directory under which relative record names are read; it does not apply to
  full paths or URLs), \code{network} and \code{legacy} are used.
\end{params}
\returns{A \code{Header} (\cref{tab:rust-header}).}
\raises{As \code{read\_header}.}
\begin{lstlisting}[language=Rust,style=example]
// A copy of the header 100s.hea in the directory "in".
std::fs::create_dir_all("in")?;
std::fs::copy("100s.hea", "in/100s.hea")?;
let opt = ReadOptions { input_dir: Some("in".into()), ..Default::default() };
let hdr = read_header_with("100s", &opt)?; // reads in/100s.hea
println!("{} {}", hdr.nsig, hdr.fs); // 2 360
println!("{}", read_header("100s", false).is_ok()); // true: the input directory is restored
\end{lstlisting}

\apientry{read\_annotations}{Read an annotation file}
\begin{lstlisting}[language=Rust,style=signature]
pub fn read_annotations(
    record: impl AsRef<[u8]>,
    annotator: impl AsRef<[u8]>,
    from: Option<TimeSpec>,
    to: Option<TimeSpec>,
    fs: Option<f64>,
    legacy: bool,
) -> Result<Annotations, Error>
\end{lstlisting}
Reads the annotations of an annotation file, in file order, as parallel vectors. The
sampling frequency, needed to convert times in seconds, is that of the record's header,
or, for a record without a header or whose header gives none, the argument \code{fs}, as
the option \opt{-F} of \code{rdann}.
\begin{params}
\item[record] Record name, as for \code{read\_record}: a full path (absolute or
  \code{\textasciitilde}), used as given; a relative path, read from the input directory
  if one is set, otherwise from the current directory; or a URL, read only if network
  access is on Its header, if there is one, gives the sampling frequency.
\item[annotator] Annotator name: the suffix of the annotation file. \code{"atr"} reads the
  file \file{RECORD.atr}.
\item[from] The first time to read: a sample number or a time string. \code{None}: the
  beginning of the file.
\item[to] The time at which to stop: annotations at or after it are not read.
  \code{None}: the end of the file.
\item[fs] The sampling frequency, in Hz ($> 0$), to use when the record has no header or
  its header gives none; it is ignored otherwise. \code{None}: none.
\item[legacy] \code{true} to behave as the original WFDB package (which assumes 250\,Hz
  for a record without a frequency); normally \code{false}.
\end{params}
\returns{An \code{Annotations} (\cref{tab:rust-annotations}). Its \code{fs} is the
frequency in which the times are counted: the one written in the annotation file, if
there is one, or else the sampling frequency; it is 0 if none of them is known.}
\raises{An \code{Error} if \code{fs} is given and is not $> 0$ (``fs must be > 0''); if
the annotation file cannot be opened (``can't read annotator ANN for record REC''); if a
time string needs the sampling frequency and none is known (``sampling frequency of
record REC is not specified''); if \code{from} cannot be reached (``can't seek in
annotator ANN''); or if the file ends early or cannot be read before \code{to} is reached
(``can't read annotation file ANN of record REC (getann status S)'', where S is the
status that \code{getann} returned, such as $-3$ for an unexpected end of file).}
\begin{lstlisting}[language=Rust,style=example]
let ann = read_annotations("100s", "atr", Some("0:10".into()), Some("0:30".into()), None, false)?;
println!("{} {}", ann.len(), ann.fs); // 24 360
println!("{} {}", ann.samples[0], s(&ann.symbols[0])); // 3862 N

// The whole file: the first annotation is a rhythm change, with an aux string.
let ann = read_annotations("100s", "atr", None, None, None, false)?;
let aux = ann.aux[0].as_deref().unwrap_or(b"");
println!("{} {} {}", s(&ann.symbols[0]), ann.codes[0] == RHYTHM, s(aux)); // + true (N
\end{lstlisting}
For a record without a header, \code{fs} gives the frequency in which times in seconds
are converted (the library reports on the standard error that it cannot open the
header):
\begin{lstlisting}[language=Rust,style=example]
let beats = Annotations {
    samples: vec![100, 460, 820],
    symbols: vec!["N".into(), "N".into(), "N".into()],
    ..Default::default()
};
write_annotations("beats", "qrs", &beats, 0.0, None)?; // no header
let ann = read_annotations("beats", "qrs", None, Some("0:02".into()), Some(360.0), false)?;
println!("{} {}", ann.len(), ann.fs); // 2 360
\end{lstlisting}

\apientry{read\_annotations\_with}{Read an annotation file, with options}
\begin{lstlisting}[language=Rust,style=signature]
pub fn read_annotations_with(
    record: impl AsRef<[u8]>,
    annotator: impl AsRef<[u8]>,
    fs: Option<f64>,
    opt: &ReadOptions,
) -> Result<Annotations, Error>
\end{lstlisting}
As \code{read\_annotations}, with the options \code{from}, \code{to}, \code{input\_dir},
\code{network} and \code{legacy} of \code{opt} (the last three for this call only). The
other fields of \code{opt} are not used.
\begin{params}
\item[record] Record name, as for \code{read\_annotations}.
\item[annotator] Annotator name, as for \code{read\_annotations}.
\item[fs] The sampling frequency to use when the record has no header or its header
  gives none, as for \code{read\_annotations}.
\item[opt] The options (\cref{tab:rust-readoptions}); only \code{from}, \code{to},
  \code{input\_dir} (optional: the input directory under which relative record names are
  read; it does not apply to full paths or URLs), \code{network} and \code{legacy} are
  used.
\end{params}
\returns{An \code{Annotations} (\cref{tab:rust-annotations}).}
\raises{As \code{read\_annotations}.}
\begin{lstlisting}[language=Rust,style=example]
// Seconds 10 to 30 of the annotations of 100s, read from the directory "in".
std::fs::create_dir_all("in")?;
for f in ["100s.hea", "100s.atr"] {
    std::fs::copy(f, format!("in/{f}"))?;
}
let opt = ReadOptions {
    from: Some("0:10".into()),
    to: Some("0:30".into()),
    input_dir: Some("in".into()),
    ..Default::default()
};
let ann = read_annotations_with("100s", "atr", None, &opt)?;
println!("{} {}", ann.len(), ann.samples[0]); // 24 3862
\end{lstlisting}

% ------------------------------------------------------------------------------------------
\subsection{Writing records}
\label{sec:rust-writing}

\apientry{write\_record}{Write a header and a signal file}
\begin{lstlisting}[language=Rust,style=signature]
pub fn write_record(
    record: impl AsRef<[u8]>,
    samples: &[i32],
    nsig: usize,
    fs: f64,
    opt: &WriteOptions,
) -> Result<(), Error>
\end{lstlisting}
Writes samples in ADC units as a new record: the header \file{RECORD.hea} and one signal
file, \file{RECORD.dat}, which holds all the signals (so they share one format). Existing
files of these names are overwritten. The sampling frequency, the gain and ADC resolution
of each signal, and the units of each calibrated signal must be given: there are no
defaults. Every check except those of the writing itself comes before any file is
created; whatever happens, the output signals are closed before the function returns,
and the getvec mode and the time resolution of output annotators are left as they were.
\begin{params}
\item[record] Record name (not empty), without the \file{.hea} extension: a plain name
  (\code{"rec"}) or a relative path (\code{"data/rec"}), created in the output directory
  if one is set, otherwise in the current directory, missing directories included; or a
  full path, absolute or in the home directory (\code{"\textasciitilde/Documents/rec"}),
  used as given. With a path, the files are created in that directory, and the header
  names the record and its signal file without the path (\code{rec}, \code{rec.dat}), so
  that the record opens under any path that leads to it (in legacy mode, as in the
  original, the header gives the name with its path). A URL cannot be written.
\item[samples] The samples, in ADC units, frame by frame: \code{nsig} values per frame,
  or, for a multi-frequency record, the sum of \code{opt.samps\_per\_frame} (the samples
  of signal 0, then those of signal 1, and so on). \code{WFDB\_INVALID\_SAMPLE} marks
  invalid samples. The number of samples must be a multiple of the frame size.
\item[nsig] The number of signals ($> 0$).
\item[fs] The sampling frequency, in Hz: a finite number $> 0$ (for a multi-frequency
  record, the frame frequency).
\item[opt] The options (\cref{tab:rust-writeoptions}): the format, the signal
  specifications (\code{gains}, \code{adcres} and \code{units} are required), the start
  time, the comments and metadata, the directory \code{dir} (which must exist) in which to
  create the files for this call only, \code{record} being then a plain name, and
  \code{legacy}.
\end{params}
\returns{\code{Ok(())}.}
\raises{An \code{Error} if a \code{samps\_per\_frame} value is less than 1
(``write\_record: samps\_per\_frame must hold integers >= 1'', checked first); with
\code{opt.dir}, if the directory does not exist (``directory D does not exist'') or the
record name contains one (``record name must not contain a directory when dir is
given''); if the record name is empty (``write\_record: the record name must not be
empty''); if \code{nsig} is 0 (``write\_record: at least one signal is required''); if a
per-signal field of \code{opt} has more than one value but not one per signal
(``write\_record: OPTION must have one value per signal''); if \code{fs} is not $> 0$ or
not finite (``write\_record: the sampling frequency must be > 0'', ``\ldots must be
finite''); if a gain is not finite (``write\_record: gains must be finite''); if an entry
of \code{opt.metadata} is not valid (``write\_record: metadata key 'K' is also given in
comments'', for example); if the gains, the ADC resolutions or the units of a calibrated
signal are missing (``write\_record: gains must be specified'', ``write\_record: ADC
resolution(s) must be specified'', ``write\_record: units must be specified''); if
\code{samples.len()} is not a multiple of the frame size (``write\_record: the number of
samples is not a multiple of the frame size''); if \code{opt.base\_time} is not a valid
time (``write\_record: incorrect base time 'T'''); if the files cannot be created
(``can't create signal file for record REC'', ``can't create header for record REC''); or
if they cannot be written completely, for example on a full disk (``write error on record
REC'').}
\begin{lstlisting}[language=Rust,style=example]
// A ramp of 1000 samples of one signal, in millivolts, at 100 Hz: the record "ramp"
// (ramp.hea and ramp.dat).
let x: Vec<i32> = (0..1000).collect();
let opt = WriteOptions {
    gains: vec![200.0],
    adcres: vec![12],
    units: vec![Some("mV".into())],
    names: vec!["ramp".into()],
    base_time: Some("14:00:03 21/06/2024".into()),
    comments: vec!["a ramp".into()],
    metadata: vec![("Generator".into(), "loop".into())],
    ..Default::default()
};
write_record("ramp", &x, 1, 100.0, &opt)?;
let rec = read_record("ramp", &ReadOptions::default())?;
println!("{} {} {}", rec.nsamples(), rec.base_time.as_deref().unwrap_or(""), rec.physical(10, 0)?);
// 1000 [14:00:03 21/06/2024] 0.05
for c in &read_header("ramp", false)?.comments {
    println!("[{}]", s(c));
}
// [a ramp]
// [ Generator: loop]
\end{lstlisting}
A multi-frequency record: signal \code{fast} has 2 samples per frame, signal
\code{slow} one; each frame holds the two samples of \code{fast} and then the sample of
\code{slow}:
\begin{lstlisting}[language=Rust,style=example]
let mut frames = Vec::new();
for i in 0..100 {
    // 1 second at 100 frames per second
    frames.push(2 * i); // fast: 200 samples per second
    frames.push(2 * i + 1);
    frames.push(500); // slow: 100 samples per second
}
let opt = WriteOptions {
    gains: vec![200.0],
    adcres: vec![12],
    units: vec![Some("mV".into())],
    names: vec!["fast".into(), "slow".into()],
    samps_per_frame: vec![2, 1],
    ..Default::default()
};
write_record("mf", &frames, 2, 100.0, &opt)?;
let hi = ReadOptions { highres: true, ..Default::default() };
let rec = read_record("mf", &hi)?;
println!("{} {} {} {} {}", rec.fs, rec.at(0, 0), rec.at(0, 1), rec.at(1, 0), rec.at(1, 1));
// 200 0 500 1 500
\end{lstlisting}

\apientry{write\_record\_physical}{Write a record from samples in physical units}
\begin{lstlisting}[language=Rust,style=signature]
pub fn write_record_physical(
    record: impl AsRef<[u8]>,
    samples: &[f64],
    nsig: usize,
    fs: f64,
    opt: &WriteOptions,
) -> Result<(), Error>
\end{lstlisting}
As \code{write\_record}, with the samples in physical units. Each sample $v$ is stored as
$\mathrm{round}(v \cdot \mathit{gain}) + \mathit{baseline}$ (halves rounded away from
zero), and NaN is stored as an invalid sample. Every signal must have a gain other than
0.
\begin{params}
\item[record] Record name, as for \code{write\_record}.
\item[samples] The samples, in physical units, frame by frame (as for
  \code{write\_record}).
\item[nsig] The number of signals ($> 0$).
\item[fs] The sampling frequency, in Hz: a finite number $> 0$.
\item[opt] The options, as for \code{write\_record} (\cref{tab:rust-writeoptions}).
\end{params}
\returns{\code{Ok(())}.}
\raises{As \code{write\_record}; also an \code{Error} if a gain is 0 (``write\_record:
signal N is uncalibrated''), or if a sample is infinite or its ADC value is outside the
range of \code{i32} (``write\_record: sample I is infinite or beyond the range of 32-bit
ADC units''). Nothing is written then.}
\begin{lstlisting}[language=Rust,style=example]
// Two sine waves, in millivolts, at 360 Hz, frame by frame: the record "sine".
let fs = 360.0;
let mut x = Vec::new();
for i in 0..3600 {
    let t = 2.0 * std::f64::consts::PI * i as f64 / fs;
    x.push(t.sin()); // signal 0
    x.push(0.5 * t.cos()); // signal 1
}
let opt = WriteOptions {
    gains: vec![1000.0],
    adcres: vec![16],
    units: vec![Some("mV".into())],
    names: vec!["A".into(), "B".into()],
    ..Default::default()
};
write_record_physical("sine", &x, 2, fs, &opt)?;
let rec = read_record("sine", &ReadOptions::default())?;
let mut err: f64 = 0.0;
for (i, v) in x.iter().enumerate() {
    err = err.max((rec.physical(i / 2, i % 2)? - v).abs());
}
println!("{} {}", rec.nsamples(), err <= 0.0005); // 3600 true: rounded to 1/1000 mV
\end{lstlisting}

\apientry{write\_annotations}{Write an annotation file}
\begin{lstlisting}[language=Rust,style=signature]
pub fn write_annotations(
    record: impl AsRef<[u8]>,
    annotator: impl AsRef<[u8]>,
    ann: &Annotations,
    fs: f64,
    dir: Option<&Path>,
) -> Result<(), Error>
\end{lstlisting}
Writes an annotation file, \file{RECORD.ANNOTATOR}, in the MIT format; an existing file of
that name is overwritten. The record need not have a header. The type of each annotation
is its \code{symbols} entry (\code{b"N"}, \code{b"V"}, \ldots; an unknown mnemonic gives
\code{NOTQRS}), or, if \code{symbols} is empty or the entry is empty, its \code{codes}
entry (\cref{chap:formats}). Both may be given (as \code{read\_annotations} returns them)
if they agree: the code of each non-empty symbol (\code{strann}) is the code given for
that annotation. \code{subtypes}, \code{chans} and \code{nums} are stored in one byte
each, as by \code{putann} (other values wrap around: the subtype and number to $-128$ to
127, the channel to 0 to 255), and \code{aux} texts are cut to 255 bytes. If the
annotations are not in time order, the file is sorted when it is written, unless the
environment variable \code{WFDBANNSORT} is \code{0} (\code{write\_annotations\_with}
decides this for one call). The arguments are checked before the file is created;
whatever happens, the annotator is closed before the function returns, and the getvec
mode and the time resolution of output annotators (\code{getafreq()}) are left as they
were (\code{fs} applies to this file only).
\begin{params}
\item[record] Record name, without the \file{.hea} extension (a trailing \code{".hea"} is
  also accepted): a plain name or a relative path (\code{"data/100s"}), created in the
  output directory if one is set, otherwise in the current directory; or a full path,
  absolute or in the home directory (\code{"\textasciitilde/Documents/100s"}), used as
  given. The file is \file{RECORD.ANNOTATOR} at that place (\file{data/100s.qrs}). A URL
  cannot be written.
\item[annotator] Annotator name: the suffix of the file. \code{"qrs"} writes
  \file{RECORD.qrs}.
\item[ann] The annotations (\cref{tab:rust-annotations}): \code{samples}, and
  \code{symbols} or \code{codes}; \code{subtypes}, \code{chans}, \code{nums} and
  \code{aux} are optional. Each of these vectors is empty (\code{subtypes}, \code{chans}
  and \code{nums} are then 0, and there is no aux) or has one entry per annotation.
  \code{ann.record}, \code{ann.annotator} and \code{ann.fs} are not used.
\item[fs] 0 (none) or a finite frequency $> 0$, in which the times are counted, recorded
  in the file; it applies to this file only.
\item[dir] The directory (which must exist) in which to create the file, for this call
  only. \code{record} must then be a plain name, without a directory: the file is
  \file{DIR/RECORD.ANNOTATOR}. \code{None}: the output directory.
\end{params}
\returns{\code{Ok(())}.}
\raises{An \code{Error} if a per-annotation field has another length than 0 or the number
of annotations (``write\_annotations: symbols must have one entry per annotation (N), or
none'', for example); if \code{symbols} and \code{codes} are both given and disagree
(``write\_annotations: symbols and codes disagree at annotation K'', K counted from 0);
if \code{fs} is not 0 and not a finite frequency $> 0$ (``write\_annotations: fs must be
finite and > 0 (or 0: none)''); if the file cannot be created (``can't create annotator
ANN for record REC''); if it cannot be written completely (``write error on annotator ANN
for record REC''); or, with \code{dir}, if the directory does not exist (``directory D
does not exist'') or \code{record} contains a directory (``record name must not contain a
directory when dir is given'').}
\begin{lstlisting}[language=Rust,style=example]
// Three beats and a rhythm change, by symbol, with an aux string for the rhythm.
let a = Annotations {
    samples: vec![18, 77, 370, 662],
    symbols: vec!["+".into(), "N".into(), "V".into(), "N".into()],
    aux: vec![Some("(N".into()), None, None, None],
    ..Default::default()
};
write_annotations("100s", "test", &a, 0.0, None)?;
let b = read_annotations("100s", "test", None, None, None, false)?;
let aux = b.aux[0].as_deref().unwrap_or(b"");
println!("{} {} {}", s(&b.symbols[0]), s(&b.symbols[2]), s(aux)); // + V (N

// A copy of the reference annotations, as read_annotations returns them.
let atr = read_annotations("100s", "atr", None, None, None, false)?;
write_annotations("100s", "copy", &atr, 0.0, None)?;
let c = read_annotations("100s", "copy", None, None, None, false)?;
println!("{}", c.samples == atr.samples && c.aux == atr.aux); // true
\end{lstlisting}

\apientry{write\_annotations\_with}{Write an annotation file, with options}
\begin{lstlisting}[language=Rust,style=signature]
pub fn write_annotations_with(
    record: impl AsRef<[u8]>,
    annotator: impl AsRef<[u8]>,
    ann: &Annotations,
    fs: f64,
    opt: &AnnWriteOptions,
) -> Result<(), Error>
\end{lstlisting}
As \code{write\_annotations}, with the output directory and the sorting of the file given
by \code{opt}, for this call only.
\begin{params}
\item[record] Record name, as for \code{write\_annotations}.
\item[annotator] Annotator name, as for \code{write\_annotations}.
\item[ann] The annotations, as for \code{write\_annotations}.
\item[fs] 0 or the frequency recorded in the file, as for \code{write\_annotations}.
\item[opt] The options (\cref{tab:rust-annwriteoptions}): \code{dir}, as the argument
  \code{dir} of \code{write\_annotations}, and \code{sort}: \code{Some(true)} always sorts
  annotations that are not in time order, \code{Some(false)} keeps the given order
  (without the library's advice to run \code{sortann}), and \code{None} lets
  \code{WFDBANNSORT} decide (sorting is on unless it is \code{0}). The environment is not
  changed.
\end{params}
\returns{\code{Ok(())}.}
\raises{As \code{write\_annotations}.}
\begin{lstlisting}[language=Rust,style=example]
let u = Annotations {
    samples: vec![500, 100],
    codes: vec![NORMAL, NORMAL],
    ..Default::default()
};
// Sorted when the file is written (the library reports it on the standard error).
write_annotations_with("100s", "sorted", &u, 0.0, &AnnWriteOptions::default())?;
let a = read_annotations("100s", "sorted", None, None, None, false)?;
println!("{:?}", a.samples); // [100, 500]
// Kept in the order given, in the existing directory "out".
std::fs::create_dir_all("out")?;
let opt = AnnWriteOptions { dir: Some("out".into()), sort: Some(false) };
write_annotations_with("100s", "unsorted", &u, 0.0, &opt)?;
println!("{}", std::path::Path::new("out/100s.unsorted").exists()); // true
\end{lstlisting}

% ------------------------------------------------------------------------------------------
\subsection{Checking records}
\label{sec:rust-checking}

\apientry{check\_record}{Check a record and its annotation files}
\begin{lstlisting}[language=Rust,style=signature]
pub fn check_record(record: impl AsRef<[u8]>, annotators: &[&str]) -> Result<CheckReport, Error>
\end{lstlisting}
Checks the header of a record, its signal specifications, every sample of its signal
files, each segment of a multi-segment record, and the annotation files named, and
reports each problem found. A problem in the record is never an error of the function: it
is reported, with a severity (\code{Severity::Error}: the record is damaged or does not
follow the WFDB format, and programs will read wrong data or fail;
\code{Severity::Warning}: the record can be read, but something is missing or unusual,
and may give wrong results; \code{Severity::Info}: a fact worth knowing, such as invalid
samples). The problems, their codes and their messages are the same in the seven
languages (\cref{tab:python-issue-codes}). The checks run in normal mode, even if legacy
mode is on, and the library prints no messages during the call (the previous settings
are restored afterwards, the getvec mode and the time resolution of output annotators
among them). They are similar to those of Benjamin Moody's \code{wfdbcheck}.
\begin{params}
\item[record] Record name, as for \code{read\_record}: a full path (absolute or
  \code{\textasciitilde}), used as given; a relative path, read from the input directory
  if one is set, otherwise from the current directory; or a URL, read only if network
  access is on.
\item[annotators] The annotator names whose files are checked, such as \code{\&["atr"]};
  \code{\&[]}: none.
\end{params}
\returns{A \code{CheckReport} (\cref{tab:rust-checkreport}): \code{report.ok} is
\code{true} if no problem is an error, \code{report.issues} lists the problems
(\code{Issue} values, \cref{tab:rust-issue}), and \code{report.lines()} (or
\code{Display}) gives one line per problem, or \code{no problems found}.}
\raises{An \code{Error} if \code{record} is empty (``record must not be empty''). A record
that cannot be read is reported (code \code{header\_unreadable}), not returned as an
error.}
\begin{lstlisting}[language=Rust,style=example]
let report = check_record("100s", &["atr"])?;
println!("{}", report.ok); // true: no errors
println!("{report}");
// warning: signal 0: the units are not specified
// warning: signal 1: the units are not specified

// A damaged record: one sample of a 12-bit signal is overwritten with 2048.
let x: Vec<i32> = (0..100).collect();
let opt = WriteOptions {
    gains: vec![100.0],
    adcres: vec![12],
    units: vec![Some("mV".into())],
    names: vec!["ramp".into()],
    ..Default::default()
};
write_record("damaged", &x, 1, 100.0, &opt)?;
let mut dat = std::fs::read("damaged.dat")?;
dat[20] = 0x00; // sample 10 (2 bytes per sample in format 16):
dat[21] = 0x08; // 0x0800 = 2048, little-endian
std::fs::write("damaged.dat", &dat)?;
let report = check_record("damaged", &[])?;
println!("{}", report.ok); // false
for i in &report.issues {
    println!("{} {} {}", i.severity, i.code, i.r#where);
}
// error sample_range signal 0
// error checksum signal 0
println!("{}", report.lines()[0]);
// error: signal 0: 1 sample is outside the ADC range (-2048 to 2047): sample 10 (value 2048)
\end{lstlisting}

\apientry{check\_record\_with}{Check a record, with options}
\begin{lstlisting}[language=Rust,style=signature]
pub fn check_record_with(
    record: impl AsRef<[u8]>,
    annotators: &[&str],
    opt: &ReadOptions,
) -> Result<CheckReport, Error>
\end{lstlisting}
As \code{check\_record}, with the options \code{input\_dir} and \code{network} of
\code{opt}, for this call only. The other fields of \code{opt} are not used.
\begin{params}
\item[record] Record name, as for \code{check\_record}.
\item[annotators] The annotator names whose files are checked, as for
  \code{check\_record}.
\item[opt] The options (\cref{tab:rust-readoptions}); only \code{input\_dir} (optional:
  the input directory under which relative record names are read; it does not apply to
  full paths or URLs) and \code{network} are used.
\end{params}
\returns{A \code{CheckReport} (\cref{tab:rust-checkreport}).}
\raises{As \code{check\_record}.}
\begin{lstlisting}[language=Rust,style=example]
std::fs::create_dir_all("in")?;
for f in ["100s.hea", "100s.dat", "100s.atr"] {
    std::fs::copy(f, format!("in/{f}"))?;
}
let opt = ReadOptions { input_dir: Some("in".into()), ..Default::default() };
let report = check_record_with("100s", &["atr"], &opt)?;
println!("{} {}", report.ok, report.issues.len()); // true 2
\end{lstlisting}

% ------------------------------------------------------------------------------------------
\subsection{Header metadata}
\label{sec:rust-metadata}

These functions read and change the metadata of a header: its info strings of the form
\code{\# Key: value}, such as \code{\# Age: 63} (\cref{chap:metadata}). In short: keys are
compared without regard to letter case (A-Z and a-z only); a key must not be empty and
must not contain \code{:}, \code{\#}, CR, LF or NUL, and a value must not contain CR, LF
or NUL; spaces, tabs, VT, FF and CR at either end of keys and values are removed; keys and
values are bytes (arguments \code{impl AsRef<[u8]>}, results \code{Vec<u8>}). A change
rewrites the header in one step, through a temporary file in the directory of the
destination, which then replaces it, keeping its permissions (on Unix) and every other
byte of it; the temporary file is removed on any failure. All the arguments are checked
first, so an error leaves the header unchanged, and nothing is written if nothing changes
(unless \code{dir} is given). Only the \file{.hea} file is read (not a \file{.info}
file); for a multi-segment record, it is the master header, unless a segment is named.
The header is opened as the library opens a header (\code{wfdb\_open}), so in legacy mode
it is searched for in the WFDB path. The entries are returned as a \code{Metadata}:
\begin{lstlisting}[language=Rust]
pub type Metadata = Vec<(Vec<u8>, Vec<u8>)>;
\end{lstlisting}
Each function has a form whose name ends in \code{\_with}, which takes a
\code{MetadataOptions} (\cref{tab:rust-metadataoptions}): the input directory and network
access for reading the header, and, for the functions that change it, the directory
\code{dir} in which to write it. These functions do not call \code{wfdbquit()}. Metadata
often holds information about patients: WFDB-x stores it as plain text, exactly as given,
and does not de-identify it.

\begin{phibox}
A header can hold identifying data about a patient (a name, a date of birth, a medical
record number). The metadata functions read and write such data as given. Records that
are shared must be de-identified first.
\end{phibox}

\apientry{read\_metadata}{Read the metadata of a record}
\begin{lstlisting}[language=Rust,style=signature]
pub fn read_metadata(record: impl AsRef<[u8]>) -> Result<Metadata, Error>
\end{lstlisting}
Reads the metadata entries of the header of a record, in file order. An info string is an
entry if it contains a colon: the key is the text before the first colon, the value the
rest of the line, both without spaces at either end. Info strings without a colon, or
whose key is empty or contains \code{\#}, are free comments, which are not returned.
\begin{params}
\item[record] Record name, without the \file{.hea} extension (a trailing \code{".hea"} is
  also accepted): a full path, absolute (\code{"/data/mitdb/100"}) or in the home
  directory (\code{"\textasciitilde/Documents/100s"}), used as given; a relative path
  (\code{"100s"}, \code{"data/100s"}), read from the input directory if one is set,
  otherwise from the current directory; or a URL
  (\code{"https://physionet.org/files/mitdb/1.0.0/100"}), read only if network access is
  on.
\end{params}
\returns{The entries, as (key, value) pairs, in file order. A key that appears more than
once (ignoring letter case) is returned once, with its first value, with a warning
(\cref{sec:rust-warnings}).}
\raises{An \code{Error} if the header cannot be found or read (``can't read header for
record REC''), or if \code{record} is empty (``record name must not be empty'') or
\code{"-"} (``the metadata of a header read from the standard input cannot be read or
changed'').}
\begin{lstlisting}[language=Rust,style=example]
let opt = WriteOptions {
    gains: vec![200.0],
    adcres: vec![12],
    units: vec![Some("mV".into())],
    names: vec!["ECG".into()],
    comments: vec!["recorded at rest".into()],
    metadata: vec![("Age".into(), "63".into()), ("Sex".into(), "F".into())],
    ..Default::default()
};
write_record("pat", &[0; 360], 1, 360.0, &opt)?;
for (key, value) in read_metadata("pat")? {
    println!("{} = {}", s(&key), s(&value));
}
// Age = 63
// Sex = F
for c in read_header("pat", false)?.comments {
    println!("[{}]", s(&c));
}
// [recorded at rest]: the free comment is not an entry
// [ Age: 63]
// [ Sex: F]
\end{lstlisting}

\apientry{read\_metadata\_with}{Read the metadata of a record, with options}
\begin{lstlisting}[language=Rust,style=signature]
pub fn read_metadata_with(
    record: impl AsRef<[u8]>,
    opt: &MetadataOptions,
) -> Result<Metadata, Error>
\end{lstlisting}
As \code{read\_metadata}, with the input directory and network access of \code{opt}, for
this call only. \code{opt.dir} is not used.
\begin{params}
\item[record] Record name, as for \code{read\_metadata}.
\item[opt] The options (\cref{tab:rust-metadataoptions}); only \code{input\_dir} and
  \code{network} are used.
\end{params}
\returns{The entries, as for \code{read\_metadata}.}
\raises{As \code{read\_metadata}.}
\begin{lstlisting}[language=Rust,style=example]
std::fs::create_dir_all("in")?;
std::fs::copy("100s.hea", "in/100s.hea")?;
let opt = MetadataOptions { input_dir: Some("in".into()), ..Default::default() };
let md = read_metadata_with("100s", &opt)?; // reads in/100s.hea
println!("{} {}", md.len(), s(&md[0].1)); // 1 0
\end{lstlisting}

\apientry{get\_metadata}{Return the value of a metadata entry}
\begin{lstlisting}[language=Rust,style=signature]
pub fn get_metadata(
    record: impl AsRef<[u8]>,
    key: impl AsRef<[u8]>,
) -> Result<Option<Vec<u8>>, Error>
\end{lstlisting}
Returns the value of one metadata entry of the header of a record.
\begin{params}
\item[record] Record name (as for \code{read\_metadata}).
\item[key] The key, such as \code{"Age"}. Letter case is ignored, and spaces at either end
  are removed; it must be a valid key (see \code{set\_metadata}), which is checked before
  the header is read.
\end{params}
\returns{\code{Some(value)} (empty for an entry with an empty value), or \code{None} if the
header has no entry for \code{key}. If it has more than one, the first is returned, with
the warning of \code{read\_metadata}.}
\raises{As \code{read\_metadata}; also an \code{Error} if \code{key} is not valid
(``metadata key 'a\#b' must not contain '\#''', for example).}
\begin{lstlisting}[language=Rust,style=example]
let opt = WriteOptions {
    gains: vec![200.0],
    adcres: vec![12],
    units: vec![Some("mV".into())],
    metadata: vec![("Age".into(), "63".into()), ("Notes".into(), "".into())],
    ..Default::default()
};
write_record("pat", &[0; 360], 1, 360.0, &opt)?;
println!("{:?}", get_metadata("pat", "age")?.map(|v| s(&v))); // Some("63")
println!("{:?}", get_metadata("pat", " NOTES ")?.map(|v| s(&v))); // Some(""): an empty value
println!("{:?}", get_metadata("pat", "Weight")?); // None: no such entry
\end{lstlisting}

\apientry{get\_metadata\_with}{Return the value of a metadata entry, with options}
\begin{lstlisting}[language=Rust,style=signature]
pub fn get_metadata_with(
    record: impl AsRef<[u8]>,
    key: impl AsRef<[u8]>,
    opt: &MetadataOptions,
) -> Result<Option<Vec<u8>>, Error>
\end{lstlisting}
As \code{get\_metadata}, with the input directory and network access of \code{opt}, for
this call only. \code{opt.dir} is not used.
\begin{params}
\item[record] Record name (as for \code{read\_metadata}).
\item[key] The key (as for \code{get\_metadata}).
\item[opt] The options (\cref{tab:rust-metadataoptions}); only \code{input\_dir} and
  \code{network} are used.
\end{params}
\returns{As \code{get\_metadata}.}
\raises{As \code{get\_metadata}.}
\begin{lstlisting}[language=Rust,style=example]
let opt = WriteOptions {
    gains: vec![200.0],
    adcres: vec![12],
    units: vec![Some("mV".into())],
    metadata: vec![("Age".into(), "63".into())],
    ..Default::default()
};
write_record("pat", &[0; 360], 1, 360.0, &opt)?;
std::fs::create_dir_all("in")?;
std::fs::copy("pat.hea", "in/pat.hea")?;
let opt = MetadataOptions { input_dir: Some("in".into()), ..Default::default() };
let age = get_metadata_with("pat", "Age", &opt)?.unwrap_or_default();
println!("{}", s(&age)); // 63
\end{lstlisting}

\apientry{set\_metadata}{Set the value of a metadata entry}
\begin{lstlisting}[language=Rust,style=signature]
pub fn set_metadata(
    record: impl AsRef<[u8]>,
    key: impl AsRef<[u8]>,
    value: impl AsRef<[u8]>,
) -> Result<(), Error>
\end{lstlisting}
Sets the value of a metadata entry in the header of a record. If the header has an entry
for \code{key} (ignoring letter case), its line is replaced by \code{\# Key: value}, at
the same position and with its line ending, with the spelling of \code{key} given here
(even if the value is the same). Otherwise the entry is added after the last info string
of the header (or, if it has none, after its last line that is not blank), with the line
ending of the record line. Every other byte of the header is kept.
\begin{params}
\item[record] Record name (as for \code{read\_metadata}). The header is changed where it
  was read, unless \code{dir} is given (\code{set\_metadata\_with}); a header read from a
  URL can only be written elsewhere, with \code{set\_metadata\_with} and \code{dir}.
\item[key] The key, such as \code{"Age"} or \code{"Date of birth"}. Spaces at either end
  are removed; it must not then be empty, and must not contain \code{:}, \code{\#}, CR,
  LF or NUL.
\item[value] The value, which may be a sentence. Spaces at either end are removed; it
  must not contain CR, LF or NUL, and may be empty (it is then written \code{\# Key:}).
\end{params}
\returns{\code{Ok(())}.}
\raises{An \code{Error} if \code{key} or \code{value} is not valid (``metadata key must not
be empty'', ``metadata key 'KEY' must not contain C'', ``value of metadata key 'KEY' must
not contain C'', where C is \code{':'}, \code{'\#'}, ``a line break'' or ``a NUL
character''); if \code{record} is empty or \code{"-"}; if the header cannot be read
(``can't read header for record REC''); if it has more than one entry for \code{key}
(``header of record REC has more than one entry for key KEY (lines L1, L2); remove them
first''); if it has no record line (``header of record REC has no record line''); if it
was read from a URL (``header of record REC was read from a URL, and cannot be changed
(give dir to write the changed header to a local directory)''); or if the changed header
cannot be written (``can't write header for record REC'').}
\begin{lstlisting}[language=Rust,style=example]
let opt = WriteOptions {
    gains: vec![200.0],
    adcres: vec![12],
    units: vec![Some("mV".into())],
    metadata: vec![("Age".into(), "63".into())],
    ..Default::default()
};
write_record("pat", &[0; 360], 1, 360.0, &opt)?;
set_metadata("pat", "AGE", "64")?; // replaces the line: "# AGE: 64"
set_metadata("pat", "Date of birth", "1961-02-03")?; // adds "# Date of birth: 1961-02-03"
for (key, value) in read_metadata("pat")? {
    println!("{} = {}", s(&key), s(&value));
}
// AGE = 64
// Date of birth = 1961-02-03
if let Err(e) = set_metadata("pat", "a:b", "1") {
    println!("{e}"); // metadata key 'a:b' must not contain ':'
}
\end{lstlisting}

\apientry{set\_metadata\_with}{Set the value of a metadata entry, with options}
\begin{lstlisting}[language=Rust,style=signature]
pub fn set_metadata_with(
    record: impl AsRef<[u8]>,
    key: impl AsRef<[u8]>,
    value: impl AsRef<[u8]>,
    opt: &MetadataOptions,
) -> Result<(), Error>
\end{lstlisting}
As \code{set\_metadata}, with the input directory and network access for reading the
header, and the directory \code{dir} in which to write the changed header, given by
\code{opt}, for this call only. With \code{dir}, the header that was read is unchanged,
and the changed header is written to the file of the same name in \code{dir}, even if
nothing changes.
\begin{params}
\item[record] Record name (as for \code{set\_metadata}). A header read from a URL can be
  changed only with \code{opt.dir}.
\item[key] The key (as for \code{set\_metadata}).
\item[value] The value (as for \code{set\_metadata}).
\item[opt] The options (\cref{tab:rust-metadataoptions}): \code{input\_dir},
  \code{network} and \code{dir}.
\end{params}
\returns{\code{Ok(())}.}
\raises{As \code{set\_metadata}; also an \code{Error} if \code{opt.dir} does not exist
(``directory D does not exist'').}
\begin{lstlisting}[language=Rust,style=example]
// The header of 100s (in the current directory) is unchanged; the changed copy is
// written to the directory "copy" (which must exist).
std::fs::create_dir_all("copy")?;
let opt = MetadataOptions { dir: Some("copy".into()), ..Default::default() };
set_metadata_with("100s", "Age", "69", &opt)?;
let age = get_metadata("copy/100s", "Age")?.unwrap_or_default();
println!("{} {:?}", s(&age), get_metadata("100s", "Age")?); // 69 None
\end{lstlisting}

\apientry{remove\_metadata}{Remove a metadata entry}
\begin{lstlisting}[language=Rust,style=signature]
pub fn remove_metadata(record: impl AsRef<[u8]>, key: impl AsRef<[u8]>) -> Result<usize, Error>
\end{lstlisting}
Removes a metadata entry from the header of a record. Every entry for \code{key}
(ignoring letter case) is removed, with its line ending, so this also repairs a header
that has more than one. Every other byte is kept. Nothing is written if the header has no
entry for \code{key}.
\begin{params}
\item[record] Record name (as for \code{set\_metadata}).
\item[key] The key (as for \code{set\_metadata}).
\end{params}
\returns{The number of lines removed (0 if the header has no entry for \code{key}).}
\raises{As \code{set\_metadata}, except that more than one entry for \code{key} is not an
error.}
\begin{lstlisting}[language=Rust,style=example]
let opt = WriteOptions {
    gains: vec![200.0],
    adcres: vec![12],
    units: vec![Some("mV".into())],
    metadata: vec![("Age".into(), "63".into()), ("Sex".into(), "F".into())],
    ..Default::default()
};
write_record("pat", &[0; 360], 1, 360.0, &opt)?;
println!("{}", remove_metadata("pat", "sex")?); // 1
println!("{}", remove_metadata("pat", "Weight")?); // 0: nothing is written
println!("{}", read_metadata("pat")?.len()); // 1
\end{lstlisting}

\apientry{remove\_metadata\_with}{Remove a metadata entry, with options}
\begin{lstlisting}[language=Rust,style=signature]
pub fn remove_metadata_with(
    record: impl AsRef<[u8]>,
    key: impl AsRef<[u8]>,
    opt: &MetadataOptions,
) -> Result<usize, Error>
\end{lstlisting}
As \code{remove\_metadata}, with the options of \code{opt}, for this call only (as for
\code{set\_metadata\_with}). With \code{dir}, the header is written to \code{dir} even if
nothing is removed.
\begin{params}
\item[record] Record name (as for \code{set\_metadata}).
\item[key] The key (as for \code{set\_metadata}).
\item[opt] The options (\cref{tab:rust-metadataoptions}): \code{input\_dir},
  \code{network} and \code{dir}.
\end{params}
\returns{The number of lines removed.}
\raises{As \code{remove\_metadata}; also an \code{Error} if \code{opt.dir} does not exist.}
\begin{lstlisting}[language=Rust,style=example]
let opt = WriteOptions {
    gains: vec![200.0],
    adcres: vec![12],
    units: vec![Some("mV".into())],
    metadata: vec![("Age".into(), "63".into()), ("Sex".into(), "F".into())],
    ..Default::default()
};
write_record("pat", &[0; 360], 1, 360.0, &opt)?;
std::fs::create_dir_all("out")?;
let opt = MetadataOptions { dir: Some("out".into()), ..Default::default() };
println!("{}", remove_metadata_with("pat", "Sex", &opt)?); // 1
println!("{} {}", read_metadata("pat")?.len(), read_metadata("out/pat")?.len()); // 2 1
\end{lstlisting}

\apientry{update\_metadata}{Make several changes to the metadata}
\begin{lstlisting}[language=Rust,style=signature]
pub fn update_metadata<K: AsRef<[u8]>, V: AsRef<[u8]>>(
    record: impl AsRef<[u8]>,
    changes: &[(K, Option<V>)],
) -> Result<(), Error>
\end{lstlisting}
Makes several changes to the metadata of a record in one write. Each (key, value) of
\code{changes} is applied in order, as \code{set\_metadata} does for \code{Some(value)}, or
as \code{remove\_metadata} does for \code{None}. All the changes are checked first: if one
is not valid, none is made.
\begin{params}
\item[record] Record name (as for \code{set\_metadata}).
\item[changes] The changes, as (key, value) pairs; a value of \code{None} removes the
  entry. No two keys may differ only in letter case. A list of removals only needs the
  type of its values: \code{\&[("Dx", None::<\&str>)]}. Empty \code{changes} change
  nothing (and write nothing).
\end{params}
\returns{\code{Ok(())}.}
\raises{As \code{set\_metadata}; also an \code{Error} if two keys differ only in letter
case (``metadata key 'K' is given more than once (keys are compared without regard to
letter case)'').}
\begin{lstlisting}[language=Rust,style=example]
let opt = WriteOptions {
    gains: vec![200.0],
    adcres: vec![12],
    units: vec![Some("mV".into())],
    metadata: vec![("Age".into(), "63".into()), ("Dx".into(), "164889003".into())],
    ..Default::default()
};
write_record("pat", &[0; 360], 1, 360.0, &opt)?;
update_metadata("pat", &[("Age", Some("64")), ("Dx", None), ("Sex", Some("F"))])?;
for (key, value) in read_metadata("pat")? {
    println!("{} = {}", s(&key), s(&value));
}
// Age = 64
// Sex = F
update_metadata("pat", &[("Sex", None::<&str>)])?; // removals only: the type of the values
println!("{}", read_metadata("pat")?.len()); // 1
\end{lstlisting}

\apientry{update\_metadata\_with}{Make several changes to the metadata, with options}
\begin{lstlisting}[language=Rust,style=signature]
pub fn update_metadata_with<K: AsRef<[u8]>, V: AsRef<[u8]>>(
    record: impl AsRef<[u8]>,
    changes: &[(K, Option<V>)],
    opt: &MetadataOptions,
) -> Result<(), Error>
\end{lstlisting}
As \code{update\_metadata}, with the options of \code{opt}, for this call only (as for
\code{set\_metadata\_with}). With \code{dir}, the header is written to \code{dir} even if
\code{changes} is empty.
\begin{params}
\item[record] Record name (as for \code{set\_metadata}).
\item[changes] The changes (as for \code{update\_metadata}).
\item[opt] The options (\cref{tab:rust-metadataoptions}): \code{input\_dir},
  \code{network} and \code{dir}.
\end{params}
\returns{\code{Ok(())}.}
\raises{As \code{update\_metadata}; also an \code{Error} if \code{opt.dir} does not
exist.}
\begin{lstlisting}[language=Rust,style=example]
let opt = WriteOptions {
    gains: vec![200.0],
    adcres: vec![12],
    units: vec![Some("mV".into())],
    metadata: vec![("Age".into(), "64".into())],
    ..Default::default()
};
write_record("pat", &[0; 360], 1, 360.0, &opt)?;
// The header "pat" in the directory "in" is unchanged; the changed copy is written to
// the directory "out" (which must exist).
std::fs::create_dir_all("in")?;
std::fs::create_dir_all("out")?;
std::fs::copy("pat.hea", "in/pat.hea")?;
let opt = MetadataOptions {
    input_dir: Some("in".into()),
    dir: Some("out".into()),
    ..Default::default()
};
update_metadata_with("pat", &[("Age", Some("65"))], &opt)?;
let a = get_metadata("in/pat", "Age")?.unwrap_or_default();
let b = get_metadata("out/pat", "Age")?.unwrap_or_default();
println!("{} {}", s(&a), s(&b)); // 64 65
\end{lstlisting}

% ------------------------------------------------------------------------------------------
\subsection{The options}
\label{sec:rust-options}

The options of the high-level functions are the fields of the structures of
\cref{tab:rust-readoptions,tab:rust-writeoptions,tab:rust-annwriteoptions,tab:rust-metadataoptions},
all public, and all with a \code{Default} implementation; the times are \code{TimeSpec}
values (\cref{tab:rust-timespec}). \code{ReadOptions}, \code{WriteOptions},
\code{AnnWriteOptions} and \code{TimeSpec} are defined in \code{wfdbx::record},
\code{MetadataOptions} in \code{wfdbx::metadata}; all are available at the root of the
crate.

\begin{xltabular}{\linewidth}{@{}lL@{}}
\caption{The fields of \code{ReadOptions}, the options of \code{read\_record} (all),
  \code{read\_header\_with} (\code{input\_dir}, \code{network}, \code{legacy}),
  \code{read\_annotations\_with} (\code{from}, \code{to}, \code{input\_dir},
  \code{network}, \code{legacy}) and \code{check\_record\_with} (\code{input\_dir},
  \code{network})}
\label{tab:rust-readoptions} \\
\toprule
\tabhead Field & \tabhead Meaning (default) \\
\midrule
\endfirsthead
\toprule
\tabhead Field & \tabhead Meaning (default) \\
\midrule
\endhead
\bottomrule
\endfoot
\code{channels: Vec<String>} & the signals to return, in this order, by number
  (\code{"0"}, \code{"1"}, \ldots) or by name (\code{"V5"}, as in \code{Record::names()},
  then as \code{findsig} finds it). Default empty: all the signals \\
\code{from: Option<TimeSpec>} & the first sample: a sample number ($\ge 0$) or a time
  string. Default \code{None}: 0 \\
\code{to: Option<TimeSpec>} & the sample after the last one: a sample number or a time
  string. Default \code{None}: the end of the record \\
\code{highres: bool} & \code{true} to read a multi-frequency record at its highest
  sampling frequency, repeating the samples of the slower signals; \code{false} averages
  the samples of the faster signals to one per frame. Default \code{false} \\
\code{pad: bool} & \code{true} to replace each invalid sample by the previous valid sample
  of the same signal. Default \code{false} \\
\code{legacy: bool} & behave as the original WFDB package, for this call only. Default
  \code{false} \\
\code{input\_dir: Option<String>} & the input directory for this call only (optional;
  see \code{wfdbx\_set\_input\_dir}): a local directory, a \code{\textasciitilde} path or
  the URL of a directory, under which relative record names are read. It does not apply
  to full paths or URLs. \code{""} means the default (\code{WFDBX\_INPUT\_DIR}, or the
  current directory). Default \code{None}: the current setting \\
\code{network: Option<bool>} & \code{Some(true)} to allow opening URLs for this call
  only, \code{Some(false)} to forbid it (see \code{wfdbx\_set\_network}). Default
  \code{None}: the current setting \\
\code{checksum\_warning: bool} & \code{true} to print a warning on the standard error when
  the samples do not match the checksums of the header (\cref{sec:rust-warnings}); the
  samples are returned either way. Default \code{true} \\
\end{xltabular}

\begin{xltabular}{\linewidth}{@{}lL@{}}
\caption{The fields of \code{WriteOptions}, the options of \code{write\_record} and
  \code{write\_record\_physical}. Each per-signal vector holds one value for all the
  signals, or one value per signal (or none, for those that have a default)}
\label{tab:rust-writeoptions} \\
\toprule
\tabhead Field & \tabhead Meaning (default) \\
\midrule
\endfirsthead
\toprule
\tabhead Field & \tabhead Meaning (default) \\
\midrule
\endhead
\bottomrule
\endfoot
\code{fmt: i32} & the storage format, the same for all the signals, such as \code{16},
  \code{212}, \code{24} or \code{516} (FLAC) (\cref{chap:formats}). Default \code{16} \\
\code{gains: Vec<f64>} & the gains, in ADC units per physical unit: required, and finite.
  A gain of 0 marks a signal as uncalibrated, which only \code{write\_record} (ADC units)
  allows. Negative gains are allowed \\
\code{baselines: Vec<i32>} & the ADC values that correspond to 0 physical units.
  Default: the \code{adczeros} values \\
\code{adczeros: Vec<i32>} & the ADC output for an input of 0\,V. Default
  \code{vec![0]} \\
\code{adcres: Vec<i32>} & the ADC resolutions, in bits: required, each at least 1 \\
\code{units: Vec<Option<Vec<u8>{}>{}>} & the physical units, such as
  \code{Some("mV".into())}: required for every signal whose gain is not 0; \code{None}
  for an uncalibrated signal \\
\code{names: Vec<Vec<u8>{}>} & the signal descriptions. Default: empty descriptions \\
\code{samps\_per\_frame: Vec<i32>} & the number of samples per frame of each signal, for a
  multi-frequency record (each $\ge 1$). Default: 1 \\
\code{base\_time: Option<Vec<u8>{}>} & the start time and date, such as
  \code{"12:00:00 01/01/2020"} (see \code{setbasetime}). Default \code{None}: none \\
\code{comments: Vec<Vec<u8>{}>} & the info strings, written to the header as comment lines
  (\code{"\#"} followed by the string). Default: none \\
\code{metadata: Vec<(Vec<u8>, Vec<u8>)>} & metadata entries (\cref{chap:metadata})
  written to the header after the comments, as \code{\# Key: value} (through
  \code{putinfo}), in this order. The keys and values follow the rules of
  \code{set\_metadata}; no two keys may differ only in letter case, and none may be the
  key of an entry in \code{comments} (an error, before any file is created). Default:
  none \\
\code{legacy: bool} & behave as the original WFDB package, for this call only: the gain
  defaults to 200, and the ADC resolution to one based on the format. Default
  \code{false} \\
\code{dir: Option<PathBuf>} & the directory (which must exist) in which the files are
  created, for this call only; the record name must then be a plain name, without a
  directory, and the record opens from that directory under its plain name. Default
  \code{None} (or an empty path): the output directory \\
\end{xltabular}

\begin{table}[htbp]
\centering
\caption{The fields of \code{AnnWriteOptions}, the options of
  \code{write\_annotations\_with}}
\label{tab:rust-annwriteoptions}
\begin{tabularx}{\linewidth}{@{}lL@{}}
\toprule
\tabhead Field & \tabhead Meaning (default) \\
\midrule
\code{dir: Option<PathBuf>} & the directory (which must exist) in which to create the
  file, for this call only; the record name must then be a plain name, without a
  directory: the file is \file{DIR/RECORD.ANNOTATOR}. Default \code{None}: the output
  directory \\
\code{sort: Option<bool>} & \code{Some(true)} always sorts annotations that are not in
  time order, \code{Some(false)} keeps the given order, for this call only. Default
  \code{None}: \code{WFDBANNSORT} decides (sorting is on unless it is \code{0}) \\
\bottomrule
\end{tabularx}
\end{table}

\begin{table}[htbp]
\centering
\caption{The fields of \code{MetadataOptions}, the options of the metadata functions whose
  names end in \code{\_with}}
\label{tab:rust-metadataoptions}
\begin{tabularx}{\linewidth}{@{}lL@{}}
\toprule
\tabhead Field & \tabhead Meaning (default) \\
\midrule
\code{input\_dir: Option<String>} & the input directory for reading the header, for this
  call only (optional), as in \code{ReadOptions}: it applies to relative record names
  only. Default \code{None}: the current setting \\
\code{network: Option<bool>} & network access for this call only, as in
  \code{ReadOptions}. Default \code{None}: the current setting \\
\code{dir: Option<PathBuf>} & for \code{set\_metadata\_with},
  \code{remove\_metadata\_with} and \code{update\_metadata\_with}: the directory (which
  must exist; a \code{\textasciitilde} path is expanded) in which the changed header is
  written, under the file name of the header that was read, instead of changing that
  header. The header is then written even if nothing changes. The signal files are not
  copied, and the output directory is not used. Not used by \code{read\_metadata\_with}
  and \code{get\_metadata\_with}. Default \code{None} (or an empty path): the header is
  changed where it was read \\
\bottomrule
\end{tabularx}
\end{table}

\begin{xltabular}{\linewidth}{@{}lL@{}}
\caption{\code{TimeSpec}, a sample number or a time string}
\label{tab:rust-timespec} \\
\toprule
\tabhead Item & \tabhead Meaning \\
\midrule
\endfirsthead
\toprule
\tabhead Item & \tabhead Meaning \\
\midrule
\endhead
\bottomrule
\endfoot
\code{TimeSpec::Sample(i64)} & a sample number \\
\code{TimeSpec::Str(Vec<u8>)} & a time string, converted with \code{strtim} for the open
  record \\
\code{From<i64>}, \code{From<i32>}, \code{From<u32>}, \code{From<usize>} & the sample
  number: \code{3600.into()} \\
\code{From<\&str>}, \code{From<String>}, \code{From<\&[u8]>} & the time string:
  \code{"0:10".into()} \\
\code{fn resolve(\&self, w: \&mut Wfdb) -> i64} & the sample number; a time string is
  converted with \code{strtim} for the record open in \code{w}, and a time of day gives
  the elapsed time from the start of the record (the absolute value of what
  \code{strtim} returns) \\
\code{fn needs\_freq(\&self, w: \&mut Wfdb) -> bool} & \code{true} if the time is a
  string that needs a sampling frequency to be converted (see
  \code{wfdbx\_time\_needs\_freq}) \\
\end{xltabular}

% ------------------------------------------------------------------------------------------
\subsection{The result types}
\label{sec:rust-results}

The readers return structures of the types \code{Record}, \code{Header} and
\code{Annotations}, whose fields and methods are listed in
\cref{tab:rust-record,tab:rust-header,tab:rust-annotations}; \code{check\_record} returns
a \code{CheckReport}, which holds \code{Issue} values
(\cref{tab:rust-checkreport,tab:rust-issue}). The signals of a record are described by
\code{Siginfo} structures (\cref{tab:rust-siginfo}), which are those of the WFDB library
(\cref{sec:rust-library}). Their fields are public, and may be changed. \code{Record},
\code{Header} and \code{Annotations} implement \code{Clone}, \code{Debug}, \code{Default}
and \code{PartialEq}; \code{CheckReport} and \code{Issue} implement \code{Clone},
\code{Debug}, \code{PartialEq} and \code{Eq}.

\begin{xltabular}{\linewidth}{@{}lL@{}}
\caption{The fields and methods of \code{Record}, the result of \code{read\_record}}
\label{tab:rust-record} \\
\toprule
\tabhead Member & \tabhead Contents \\
\midrule
\endfirsthead
\toprule
\tabhead Member & \tabhead Contents \\
\midrule
\endhead
\bottomrule
\endfoot
\code{name: Vec<u8>} & the record name, as given to \code{read\_record} \\
\code{fs: f64} & the sampling frequency of the returned samples, in Hz, or 0 if the header
  gives none; with \code{highres}, the highest sampling frequency of a multi-frequency
  record \\
\code{nsig: usize} & the number of signals returned \\
\code{signals: Vec<usize>} & the number, in the record, of each signal returned \\
\code{samples: Vec<i32>} & the samples in ADC units, frame by frame:
  \code{samples[i * nsig + s]} is sample \code{i} of signal \code{s}; invalid samples are
  \code{WFDB\_INVALID\_SAMPLE} \\
\code{siginfo: Vec<Siginfo>} & one per signal returned, in the order of the signals;
  \code{desc} has no white space at either end \\
\code{base\_time: Option<String>} & the start time and date as text, to the second, such
  as \code{"[14:00:03 21/06/2024]"}; \code{None} if the header gives none \\
\code{base\_time\_ms: Option<i64>} & the start time in milliseconds after midnight
  (\code{"14:00:03.500"} gives \code{50403500}); \code{Some} exactly when
  \code{base\_time} is \\
\code{base\_date: Option<[i32; 3]>} & the start date as \code{[year, month, day]};
  \code{None} if the header gives no date \\
\code{sampfrom: i64} & the number of the first sample returned \\
\code{comments: Vec<Vec<u8>{}>} & the info strings of the record: the comments of the
  header, and the lines of its \file{.info} file if it has one (see \code{getinfo}) \\
\code{checksum\_ok: bool} & \code{false} if the samples do not match the checksums of the
  header; checksums are checked only when a read reaches the end of the record, so a
  shorter read gives \code{true} \\
\code{legacy: bool} & \code{true} if the record was read in legacy mode;
  \code{physical()} then uses a gain of 200 for uncalibrated signals \\
\code{fn nsamples(\&self) -> usize} & the number of samples per signal (the number of
  frames) \\
\code{fn at(\&self, i: usize, s: usize) -> i32} & sample \code{i} of signal \code{s}, in
  ADC units; panics if \code{i} or \code{s} is out of range \\
\code{fn physical(\&self, i: usize, s: usize)} & a \code{Result<f64, Error>}: sample
  \code{i} of signal \code{s} in physical units, $(\mathit{value} - \mathit{baseline}) /
  \mathit{gain}$; NaN if the sample is invalid. An \code{Error} (``signal N of record REC
  is uncalibrated'') if the gain of the signal is 0 \\
\code{fn physical\_signal(\&self, s: usize)} & a \code{Result<Vec<f64>, Error>}: all the
  samples of signal \code{s} in physical units (the same error) \\
\code{fn names(\&self) -> Vec<Vec<u8>{}>} & the signal descriptions (empty if absent),
  without the white space that a header may put at either end of them \\
\end{xltabular}

\begin{xltabular}{\linewidth}{@{}lL@{}}
\caption{The fields of \code{Header}, the result of \code{read\_header} and
  \code{read\_header\_with}}
\label{tab:rust-header} \\
\toprule
\tabhead Field & \tabhead Contents \\
\midrule
\endfirsthead
\toprule
\tabhead Field & \tabhead Contents \\
\midrule
\endhead
\bottomrule
\endfoot
\code{name: Vec<u8>} & the record name, as given \\
\code{fs: f64} & the sampling frequency in Hz, as \code{sampfreq} gives it; 0 if the
  header gives none \\
\code{counter\_freq: f64} & the counter frequency (\code{getcfreq}) \\
\code{base\_counter: f64} & the base counter value (\code{getbasecount}) \\
\code{nsamples: i64} & the number of samples per signal, from the record line of the
  header (0 if the header does not give it) \\
\code{nsig: usize} & the number of signals \\
\code{siginfo: Vec<Siginfo>} & the signal specifications, one per signal; \code{desc} has
  no white space at either end \\
\code{base\_time}, \code{base\_time\_ms}, \code{base\_date} & the start time and date, as
  in \code{Record} \\
\code{comments: Vec<Vec<u8>{}>} & the info strings of the record, as in \code{Record} \\
\code{metadata: Vec<(Vec<u8>, Vec<u8>)>} & the metadata entries of the header, in file
  order, as \code{read\_metadata} returns them: from the \file{.hea} file only; empty,
  without an error, for an EDF record and for a header read from the standard input \\
\end{xltabular}

\begin{xltabular}{\linewidth}{@{}lL@{}}
\caption{The fields and methods of \code{Annotations}, the result of
  \code{read\_annotations} and the argument of \code{write\_annotations}: parallel
  vectors, annotation \code{i} being \code{samples[i]}, \code{codes[i]},
  \code{symbols[i]}, \ldots, in file order}
\label{tab:rust-annotations} \\
\toprule
\tabhead Member & \tabhead Contents \\
\midrule
\endfirsthead
\toprule
\tabhead Member & \tabhead Contents \\
\midrule
\endhead
\bottomrule
\endfoot
\code{record: Vec<u8>}, \code{annotator: Vec<u8>} & the record and annotator names, such
  as \code{b"100s"} and \code{b"atr"} \\
\code{fs: f64} & the frequency in which \code{samples} are counted, normally the sampling
  frequency of the record; 0 if unknown. \code{write\_annotations} does not use it \\
\code{samples: Vec<i64>} & the annotation times, in samples \\
\code{codes: Vec<i32>} & the annotation codes, such as \code{NORMAL} or \code{PVC} \\
\code{symbols: Vec<Vec<u8>{}>} & the mnemonics of the codes, as \code{annstr} returns them,
  such as \code{b"N"}, \code{b"V"} or \code{b"+"} \\
\code{subtypes: Vec<i32>}, \code{chans: Vec<i32>}, \code{nums: Vec<i32>} & the subtype,
  channel number and annotator number of each annotation \\
\code{aux: Vec<Option<Vec<u8>{}>{}>} & the auxiliary text of each annotation, without its
  length byte and up to its first NUL byte, or \code{None} where there is none \\
\code{fn len(\&self) -> usize} & the number of annotations \\
\code{fn is\_empty(\&self) -> bool} & \code{true} if there are no annotations \\
\end{xltabular}

\begin{table}[htbp]
\centering
\caption{The fields and methods of \code{CheckReport}, the result of \code{check\_record}}
\label{tab:rust-checkreport}
\begin{tabularx}{\linewidth}{@{}lL@{}}
\toprule
\tabhead Member & \tabhead Contents \\
\midrule
\code{record: Vec<u8>} & the record name, as given \\
\code{ok: bool} & \code{true} if no issue has the severity \code{Severity::Error} \\
\code{issues: Vec<Issue>} & the problems, in the order of the checks \\
\code{fn lines(\&self) -> Vec<String>} & the text form of the report: one line per issue
  (the \code{Display} form of the \code{Issue}), or the one line \code{no problems
  found} \\
\code{Display} & the lines of \code{lines()}, joined by line feeds (no final line feed) \\
\bottomrule
\end{tabularx}
\end{table}

\begin{table}[htbp]
\centering
\caption{The fields of \code{Issue}, a problem found by \code{check\_record}, and of
  \code{Severity}}
\label{tab:rust-issue}
\begin{tabularx}{\linewidth}{@{}lL@{}}
\toprule
\tabhead Member & \tabhead Contents \\
\midrule
\code{severity: Severity} & \code{Severity::Error}, \code{Severity::Warning} or
  \code{Severity::Info}; \code{as\_str()} and \code{Display} give \code{"error"},
  \code{"warning"} or \code{"info"} \\
\code{code: \&'static str} & a stable identifier of the kind of problem, such as
  \code{"checksum"} or \code{"gain\_missing"} (\cref{tab:python-issue-codes}) \\
\code{r\#where: String} & where the problem is: \code{""} (the record),
  \code{"signal N"}, \code{"segment SEG"}, \code{"segment SEG, signal N"} or
  \code{"annotator ANN"} (the field is named \code{where}, a keyword of Rust, written
  \code{r\#where}) \\
\code{message: String} & the description of the problem \\
\code{Display} & \code{SEVERITY: WHERE: MESSAGE}, or \code{SEVERITY: MESSAGE} when
  \code{where} is empty \\
\bottomrule
\end{tabularx}
\end{table}

\begin{xltabular}{\linewidth}{@{}lL@{}}
\caption{The fields of \code{Siginfo} (C's \code{WFDB\_Siginfo}), the specification of a
  signal}
\label{tab:rust-siginfo} \\
\toprule
\tabhead Field & \tabhead Contents \\
\midrule
\endfirsthead
\toprule
\tabhead Field & \tabhead Contents \\
\midrule
\endhead
\bottomrule
\endfoot
\code{fname: Option<Vec<u8>{}>} & the name of the signal file, as the header gives it \\
\code{desc: Option<Vec<u8>{}>} & the signal description, such as \code{b"MLII"} \\
\code{units: Option<Vec<u8>{}>} & the physical units, such as \code{b"mV"}; \code{None} if
  not specified \\
\code{gain: f64} & the gain, in ADC units per physical unit (0: uncalibrated) \\
\code{initval: i32} & the value of the first sample \\
\code{group: u32} & the signal group number (the signals of one group share one file) \\
\code{fmt: i32} & the storage format, such as \code{16}, \code{212} or \code{516} \\
\code{spf: i32} & the number of samples per frame (more than 1 for the faster signals of a
  multi-frequency record) \\
\code{bsize: i32} & the block size, in bytes (0 for ordinary files) \\
\code{adcres: i32} & the ADC resolution, in bits (0: not specified) \\
\code{adczero: i32} & the ADC output for an input of 0\,V \\
\code{baseline: i32} & the ADC value that corresponds to 0 physical units \\
\code{nsamp: i64} & the number of samples (0: not specified) \\
\code{cksum: i32} & the 16-bit checksum of the samples \\
\end{xltabular}

% ------------------------------------------------------------------------------------------
\subsection{Warnings and errors}
\label{sec:rust-warnings}

\paragraph{Errors.} Each high-level function returns a \code{Result}; its error type is
\code{wfdbx::record::Error}, the one public item of the high-level API that is not
available at the root of the crate. \code{e.message()} and \code{Display} give the
message, which describes the failed operation, such as ``can't read record REC''; the
type implements \code{std::error::Error}, \code{Clone}, \code{Debug} and \code{PartialEq},
so it converts to \code{Box<dyn std::error::Error>} with the \code{?} operator. There is
one error type: the messages tell the cases apart (\cref{tab:rust-errors}). Each function
checks its arguments before it creates or changes a file, so an invalid argument leaves
the files as they were. The library may also print a message of its own on the standard
error when it cannot open a file (as in \code{init: can't open header for record REC});
\code{wfdberror()} returns the last such message, and \code{wfdbquiet()} turns these
messages off (\cref{sec:rust-library}). \code{Record::at}, like the indexing of the
result vectors, panics on an index out of range.

\begin{table}[htbp]
\centering
\caption{The kinds of errors of the high-level API}
\label{tab:rust-errors}
\begin{tabularx}{\linewidth}{@{}lL@{}}
\toprule
\tabhead Kind & \tabhead Messages \\
\midrule
files & a record, header or annotation file cannot be found, read, created or written
  completely (for example on a full disk): ``can't read record REC'', ``can't read header
  for record REC'', ``can't write header for record REC'', ``write error on record REC'',
  ``directory D does not exist'', \ldots; a signal or annotation file ends before its
  header or its end-of-file marker says (``\ldots(getvec status S)'', ``\ldots(getann
  status S)'') \\
arguments & an invalid argument or option, reported before anything is written (most
  messages of \code{write\_record} and \code{write\_annotations} begin with the function's
  name); for the metadata functions, an invalid key or value, a record name that is empty
  or \code{"-"}, a key that the header has more than once, a header without a record
  line, or a header read from a URL to be changed without \code{dir} \\
missing values & the operation needs a value that the record does not give: a time in
  seconds without a sampling frequency (``sampling frequency of record REC is not
  specified''), physical units without a gain (``signal N of record REC is
  uncalibrated'') \\
\bottomrule
\end{tabularx}
\end{table}

\paragraph{Warnings.} Warnings are printed on the standard error with the library's
messages, through \code{wfdb\_error}: \code{wfdbquiet()} silences them, and
\code{wfdberror()} returns the last one. There is no callback.

\paragraph{Settings.} Whether a function returns a result or an error, the library's
settings are those it had before the call: the input and output directories, network
access and legacy mode (set for the call by the options), and, for \code{read\_record},
\code{read\_header}, \code{read\_annotations}, \code{write\_record},
\code{write\_record\_physical}, \code{write\_annotations}, their \code{\_with} forms and
\code{check\_record}, the getvec mode (\code{getgvmode()}) and the time resolution of
output annotators (\code{getafreq()}), which opening an EDF record sets (the
\code{WFDB\_HIGHRES} bit, and the record's sampling frequency when no resolution is set)
and \code{wfdbquit()} does not reset. A later \code{write\_annotations} therefore records
no frequency that an EDF record read before it left behind; its own \code{fs} applies to
its file only.

\paragraph{The checksum warning.} When \code{read\_record} reads to the end of a record
whose samples do not match the checksums of its header, it returns all the samples and
prints the warning ``WFDB-x: checksum error in record REC: the signal data do not match
the checksums in its header''. \code{checksum\_warning: false} turns the warning off;
\code{checksum\_ok} of the result is \code{false} either way.

\paragraph{The duplicate-key warning.} When a header holds two metadata entries whose keys
differ only in letter case (written by other software), \code{read\_metadata},
\code{get\_metadata} and \code{read\_header} use the first and print the warning ``WFDB-x:
header of record REC has more than one entry for key KEY (lines L1, L2); the first is
used'' (KEY spelled as in its first entry). \code{set\_metadata} and
\code{update\_metadata} then refuse to change that key (an \code{Error}), and
\code{remove\_metadata} removes all of its entries.

\begin{lstlisting}[language=Rust,style=example]
// A record whose last sample is changed after it was written: the checksum fails.
let x: Vec<i32> = (0..100).collect();
let opt = WriteOptions {
    gains: vec![100.0],
    adcres: vec![12],
    units: vec![Some("mV".into())],
    ..Default::default()
};
write_record("ck", &x, 1, 100.0, &opt)?;
let mut dat = std::fs::read("ck.dat")?;
let n = dat.len();
dat[n - 2] = 0;
dat[n - 1] = 0;
std::fs::write("ck.dat", &dat)?;
let rec = read_record("ck", &ReadOptions::default())?; // prints the checksum warning
println!("{} {}", rec.nsamples(), rec.checksum_ok); // 100 false: all the samples
let quiet = ReadOptions { checksum_warning: false, ..Default::default() };
let rec = read_record("ck", &quiet)?; // no warning

// Two entries for one key: the first is used, with a warning.
let dup = WriteOptions {
    gains: vec![0.0],
    adcres: vec![12],
    comments: vec![" Age: 63".into(), " AGE: 64".into()],
    ..Default::default()
};
write_record("dup", &[0; 10], 1, 100.0, &dup)?;
let age = get_metadata("dup", "age")?.unwrap_or_default(); // with a warning
println!("{}", s(&age)); // 63
if let Err(e) = set_metadata("dup", "Age", "65") {
    println!("{e}");
    // header of record dup has more than one entry for key Age (lines 3, 4); remove
    // them first
}
println!("{}", remove_metadata("dup", "age")?); // 2

match read_record("nosuchrecord", &ReadOptions::default()) {
    Err(e) => println!("{}", e.message()), // can't read record nosuchrecord
    Ok(_) => {}
}
\end{lstlisting}

% ==========================================================================================
\section{The WFDB library functions in Rust}
\label{sec:rust-library}

Every function of the WFDB library is a method of \code{Wfdb} and a free function of the
crate, under its C name, with the same arguments, in the same order, and the same return
values, so that programs and documentation written for the C library carry over
directly. The functions, their arguments and their results are documented once, for all
the languages, in \cref{chap:library}; this section gives the rules by which the C
declarations become Rust functions (see also \cref{sec:library-rust}). The doc comment of
each function summarizes it, and says what WFDB-x changes in it.

\subsection{Conventions}

\paragraph{The library state.} The state of the C library (the open record, its signals,
the annotators, the WFDB path, the time base, \ldots) is held by a \code{Wfdb} value
(\code{Wfdb::new()}), and each C function is a method of it with the C name:
\code{w.isigopen(...)}, \code{w.getvec(...)}. Each method is also a free function with
the same name and arguments (\code{wfdbx::isigopen(...)}), which uses the default
\code{Wfdb} value of the calling thread; a free function must not be called from inside a
closure given to \code{wfdbx::with\_default}, which borrows that value. Separate
\code{Wfdb} values share no state, unlike the C library. The functions that use no state
are free functions only (\code{wfdbldflags}, \code{wfdbcflags}, \code{wfdbx\_write\_errors},
\code{isann}, \ldots; \cref{sec:library-rust}).

\paragraph{Names and modules.} The functions, the constants (\code{WFDB\_READ},
\code{WFDB\_INVALID\_SAMPLE}, \code{WFDB\_HIGHRES}, \ldots), the structures and the
annotation codes (\code{NORMAL}, \code{PVC}, \ldots, \cref{chap:formats}) keep their C
names, and are available at the root of the crate (\code{use wfdbx::*;}). They are
defined in the modules of \cref{tab:rust-modules}. The macros of
\file{<wfdb/ecgmap.h>} are functions (\code{isann}, \code{isqrs}, \code{map1}, \ldots,
and \code{wfdb\_isann}, \code{wfdb\_isqrs}, \ldots).

\begin{table}[htbp]
\centering
\caption{The modules of the crate \code{wfdbx}}
\label{tab:rust-modules}
\begin{tabularx}{\linewidth}{@{}lL@{}}
\toprule
\tabhead Module & \tabhead Contents \\
\midrule
\code{record} & the high-level functions of reading and writing, their option and result
  types, and \code{Error} \\
\code{check} & \code{check\_record}, \code{check\_record\_with}, \code{CheckReport},
  \code{Issue}, \code{Severity} \\
\code{metadata} & the metadata functions, \code{MetadataOptions}, \code{Metadata} \\
\code{signal} & signals and headers (port of \file{lib/signal.c}) \\
\code{annot} & annotations (\file{lib/annot.c}) \\
\code{calib} & the calibration list (\file{lib/calib.c}) \\
\code{wfdbio} & file names, the input and output directories, network access, error
  messages, \code{WfdbFile} (\file{lib/wfdbio.c}) \\
\code{wfdbinit} & \code{wfdbinit}, \code{wfdbquit}, \code{wfdbflush}
  (\file{lib/wfdbinit.c}) \\
\code{ecgcodes} & the annotation codes and the functions of \file{<wfdb/ecgmap.h>} \\
\code{types} & the types, constants and structures of \file{<wfdb/wfdb.h>} \\
\code{wfdb} & the type \code{Wfdb}, \code{with\_default}, \code{OptBytes}, \code{NULL} \\
\code{cstd} & the C standard functions that the tools use, such as a \code{printf} that
  formats numbers exactly as the GNU C library does (\code{cformat!}, \code{cprintf!}) \\
\bottomrule
\end{tabularx}
\end{table}

\paragraph{Arguments.} String arguments are byte strings: they accept \code{\&str},
\code{\&[u8]} or \code{b"..."}, and a byte string ends at its first NUL byte, as a C
string does. Where C accepts \code{NULL}, they are \code{impl OptBytes}, which also
accepts \code{Option<\&[u8]>}, \code{Option<\&str>} and \code{\&Option<Vec<u8>{}>};
\code{wfdbx::NULL} is a typed \code{None}: \code{sampfreq(NULL)} returns the frequency of
the open record, and \code{getinfo(NULL)} continues with the info strings of the last
record. Where C fills an array, the caller passes a mutable slice
(\code{isigopen(record, \&mut si, n)}, \code{getvec(\&mut v)}); an empty slice stands for
C's \code{NULL} (\code{isigopen(record, \&mut [], 0)}); \code{osigopen(record, \&mut [], 0)}
closes the open output signals (unless the record name begins with \code{+}), reads the
header, opens nothing and returns 0. Where C fills a structure, the
caller passes a mutable reference (\code{getann(0, \&mut annot)}). \code{getseginfo}
takes an \code{\&mut Option<Vec<Seginfo>{}>}, which receives a copy of the segment table.

\paragraph{Strings} are returned by value, as \code{Vec<u8>}, or as
\code{Option<Vec<u8>{}>} where the C function may return \code{NULL} (\code{anndesc},
\code{getinfo}, \code{wfdbfile}, \ldots); unlike C's pointers to static buffers, they stay
valid after later calls. Text is handled as bytes.

\paragraph{Return values and errors.} The functions return what the C functions return:
an \code{i32} status (negative on errors), a number of signals, an \code{f64}, a string
or \code{None}. Errors are reported by these values and by a message on the standard
error, as in C, never by \code{Result} or panics. \code{wfdberror()} returns the last
message; \code{wfdbquiet()} turns the messages off and \code{wfdbverbose()} on again.
\code{wfdb\_error(msg)} takes a message already formatted, with the macro
\code{cformat!} (which formats as C's \code{printf}), and \code{wfdb\_error\_fmt(fmt,
args)} formats it.

\paragraph{State.} As in C, a \code{Wfdb} value keeps its state between calls: the open
record, its input signals, the output signals and the annotators remain open until
\code{wfdbquit()} closes them. \code{wfdbquit()} leaves the input and output directories
and the network setting as they are, and so, as in C, the time resolution of output
annotators (\code{getafreq()}), which opening an EDF record sets to its sampling frequency
when none is set, and the getvec mode (\code{getgvmode()}), to which it adds
\code{WFDB\_HIGHRES}. Dropping a \code{Wfdb} value flushes its output
buffers, as C's \code{exit} does, but does not close its records (no annotation end
marks, no header updates): a program calls \code{wfdbquit()} before it ends. The
high-level functions of \cref{sec:rust-api} (except the metadata functions) call
\code{wfdbquit()} on the \code{Wfdb} value they use, which closes anything opened with the
library functions on that value, and leave those two settings as they were.

\paragraph{Additions.} \code{getvecs(buffer, n)} reads up to \code{n} sample vectors at
once, with exactly the result of \code{n} calls of \code{getvec}, but decoding whole
blocks of the signal files; \code{getvecs\_status(buffer, n)} also gives the status of
the last \code{getvec} call (\cref{sec:library-additions}):
\begin{lstlisting}[language=Rust]
pub fn getvecs(out: &mut [i32], n: i64) -> i64
pub fn getvecs_status(out: &mut [i32], n: i64) -> (i64, i32)
\end{lstlisting}
\code{WfdbFile} is C's \code{FILE *} (or \code{WFDB\_FILE}): \code{wfdb\_open} and
\code{wfdb\_fopen} return an \code{Option<WfdbFile>}, whose methods (\code{getc},
\code{read}, \code{fgets}, \code{write}, \code{seek}, \code{close}, \ldots) are those of
the C standard library.

\subsection{The structures}
\label{sec:rust-structures}

The structures of the C library have the names of the C structures without the prefix
\code{WFDB\_}, and the C field names, with C's zero values as \code{Default}
(\cref{tab:rust-structures}).

\begin{xltabular}{\linewidth}{@{}lL@{}}
\caption{The C structures in Rust}
\label{tab:rust-structures} \\
\toprule
\tabhead Structure & \tabhead Fields \\
\midrule
\endfirsthead
\toprule
\tabhead Structure & \tabhead Fields \\
\midrule
\endhead
\bottomrule
\endfoot
\code{Siginfo} & \code{fname}, \code{desc}, \code{units} (\code{Option<Vec<u8>{}>}),
  \code{gain}, \code{initval}, \code{group}, \code{fmt}, \code{spf}, \code{bsize},
  \code{adcres}, \code{adczero}, \code{baseline}, \code{nsamp}, \code{cksum}
  (\cref{tab:rust-siginfo}) \\
\code{Seginfo} & \code{recname} (the name of the segment, \code{Vec<u8>}), \code{nsamp}
  (its number of samples, \code{i64}), \code{samp0} (the number of its first sample in the
  record, \code{i64}) \\
\code{Anninfo} & \code{name} (the annotator name, \code{Vec<u8>}), \code{stat}
  (\code{i32}: \code{WFDB\_READ} or \code{WFDB\_WRITE} for the MIT format,
  \code{WFDB\_AHA\_READ} or \code{WFDB\_AHA\_WRITE} for the AHA format) \\
\code{Annotation} & \code{time} (\code{i64}, in sample intervals), \code{anntyp}
  (\code{i8}), \code{subtyp} (\code{i8}), \code{chan} (\code{u8}), \code{num}
  (\code{i8}), \code{aux} (\code{Option<Vec<u8>{}>}: absent, or bytes whose first byte is
  the length of the text that follows, as in C) \\
\code{Calinfo} & \code{low}, \code{high} (the levels of the calibration pulse, in physical
  units, \code{f64}), \code{scale} (physical units per cm), \code{sigtype} and
  \code{units} (\code{Option<Vec<u8>{}>}), \code{caltype} (\code{i32}) \\
\end{xltabular}

The integer types follow the C reference platform: \code{WFDB\_Sample} and C's
\code{int} are \code{i32}, \code{WFDB\_Time}, \code{WFDB\_Date} and C's \code{long} are
\code{i64}, \code{WFDB\_Frequency} and \code{WFDB\_Gain} are \code{f64}, C's \code{char}
is \code{i8}, and \code{WFDB\_Signal}, \code{WFDB\_Group} and \code{WFDB\_Annotator} are
\code{u32}. C's wrap-around arithmetic is reproduced.

\subsection{The settings functions}
\label{sec:rust-settings}

The functions of \cref{tab:rust-settings} are additions of WFDB-x to the library: they
set where files are read and written, network access and legacy mode, for the
high-level functions, the library functions and the tools alike
(\cref{chap:files,chap:behavior}), and give the path under which a program opens a file
that it reads or creates itself, as the tools do (\code{wfdbx\_input\_path},
\code{wfdbx\_output\_path}; \cref{sec:tool-files}). Like the other library functions, each is a method of
\code{Wfdb} and a free function, except \code{wfdbx\_write\_errors} and
\code{wfdbx\_note\_write\_error} (free functions only: the count belongs to the process).
Each \code{set} function returns the previous setting, so that it can be restored. The C
function \code{setwfdb(p)} sets the input directory too (\code{p} must name at most one
directory), and \code{getwfdb()} returns it (\code{"."} if none is set); in legacy mode
they set and return the WFDB path, as in C.

\begin{xltabular}{\linewidth}{@{}lL@{}}
\caption{The settings functions (the signatures of the free functions)}
\label{tab:rust-settings} \\
\toprule
\tabhead Function & \tabhead Effect and result \\
\midrule
\endfirsthead
\toprule
\tabhead Function & \tabhead Effect and result \\
\midrule
\endhead
\bottomrule
\endfoot
\code{wfdbx\_set\_input\_dir(dir: impl OptBytes) -> Option<Vec<u8>{}>} & sets the input
  directory: a local directory, a \code{\textasciitilde} path or a URL; \code{NULL} or
  \code{""} restores the default (\code{WFDBX\_INPUT\_DIR}, or else the current
  directory). Returns the previous input directory, or \code{None} if none was set \\
\code{wfdbx\_get\_input\_dir() -> Option<Vec<u8>{}>} & returns the input directory, or
  \code{None} if none is set \\
\code{wfdbx\_set\_output\_dir(dir: impl OptBytes) -> Option<Vec<u8>{}>} & sets the output
  directory: a local directory or a \code{\textasciitilde} path (a URL is refused, with a
  message, and the setting is unchanged); \code{NULL} or \code{""} restores the default
  (\code{WFDBX\_OUTPUT\_DIR}, or else the current directory). Returns the previous output
  directory, or \code{None} \\
\code{wfdbx\_get\_output\_dir() -> Option<Vec<u8>{}>} & returns the output directory, or
  \code{None} if none is set \\
\code{wfdbx\_set\_network(on: i32) -> i32} & turns network access on ($> 0$) or off (0);
  a negative number restores the default (on if \code{WFDBX\_NETWORK} is a nonzero
  number). Returns the previous setting, 1 or 0 \\
\code{wfdbx\_get\_network() -> i32} & returns 1 if network access is on, 0 otherwise \\
\code{wfdbx\_set\_legacy(on: i32) -> i32} & turns legacy mode on (nonzero) or off (0).
  Returns the previous setting \\
\code{wfdbx\_get\_legacy() -> i32} & returns 1 if legacy mode is on, 0 otherwise \\
\code{wfdbx\_legacy\_args(argv: \&mut Vec<Vec<u8>{}>) -> i32} & if \code{argv[1]} is
  \code{"-{}-legacy"}, removes it from the arguments, turns legacy mode on and returns 1;
  otherwise returns 0. It also enables the exit check of the tools in
  \code{cstd::exit} (an exit status of 0 becomes 1 after a write error, outside legacy
  mode) \\
\code{wfdbx\_time\_needs\_freq(string: impl AsRef<[u8]>) -> i32} & returns 1 if
  \code{strtim(string)} needs a sampling frequency, 0 otherwise \\
\code{wfdbx\_write\_errors() -> i32} & returns the number of write errors since the
  program started (output files that could not be written, flushed or closed) \\
\code{wfdbx\_note\_write\_error()} & counts a write error, for a program that writes a file
  itself \\
\code{wfdbx\_input\_path(name: impl OptBytes) -> Option<Vec<u8>{}>} & returns the path
  under which a program opens a file that it reads itself: in normal mode, a relative
  name under the input directory, a \code{\textasciitilde} name expanded, an absolute name
  as given; in legacy mode, the name as given. \code{NULL}, \code{""} and \code{"-"} are
  returned as given; \code{None} for a URL, or a name that becomes one under the input
  directory \\
\code{wfdbx\_output\_path(name: impl OptBytes) -> Option<Vec<u8>{}>} & returns the path
  under which a program creates a file that it writes itself: in normal mode, a relative
  name under the output directory (whose missing directories are then created), a
  \code{\textasciitilde} name expanded, an absolute name as given; in legacy mode, any
  name under the output directory, as a plain prefix. \code{NULL}, \code{""} and
  \code{"-"} are returned as given; \code{None} for a URL \\
\end{xltabular}

\subsection{Examples}

Reading the samples and the annotations of a record, as a C program does (the record
\code{100s} is read from the current directory):

\begin{lstlisting}[language=Rust,style=example]
use wfdbx::*;

fn main() {
    let mut si = vec![Siginfo::default(); 2]; // receives one Siginfo per signal
    let nsig = isigopen("100s", &mut si, 2); // open 2 signals
    let desc = si[0].desc.clone().unwrap_or_default();
    println!("{} {} {} {}", nsig, String::from_utf8_lossy(&desc), si[0].fmt, si[0].gain);
    // 2 MLII 212 200
    println!("{}", sampfreq(NULL)); // 360: the frequency of the open record

    let mut v = [0i32; 2]; // receives one sample per signal
    isigsettime(strtim("0:30")); // skip to 30 seconds
    getvec(&mut v);
    println!("{} {} {}", v[0], v[1], aduphys(0, v[0])); // 947 968 -0.385

    let mut frames = vec![0i32; 360 * 2]; // the next second, in one call
    let (n, status) = getvecs_status(&mut frames, 360);
    println!("{n} {status}"); // 360 2

    let ai = Anninfo { name: b"atr".to_vec(), stat: WFDB_READ };
    annopen("100s", &[ai], 1);
    let mut a = Annotation::default();
    let mut count = 0;
    while getann(0, &mut a) == 0 {
        // 0: an annotation was read; -1: the end
        if count < 2 {
            let t = mstimstr(a.time);
            println!("{} {}", String::from_utf8_lossy(&t), String::from_utf8_lossy(&annstr(a.anntyp as i32)));
        }
        count += 1;
    }
    println!("{count} annotations"); // 75 annotations

    wfdbquiet(); // no messages on the standard error
    let stat = isigopen("nosuchrecord", &mut [], 0);
    print!("{} {}", stat, String::from_utf8_lossy(&wfdberror()));
    // -1 init: can't open header for record nosuchrecord
    wfdbverbose();
    wfdbquit(); // close everything
}
\end{lstlisting}

Writing a record and an annotation file, sample by sample and annotation by annotation:

\begin{lstlisting}[language=Rust,style=example]
use wfdbx::*;

fn main() {
    let out = Siginfo {
        fname: Some(b"ramp.dat".to_vec()),
        desc: Some(b"ramp".to_vec()),
        units: Some(b"mV".to_vec()),
        gain: 200.0,
        fmt: 16,
        adcres: 12,
        ..Default::default()
    };
    println!("{}", osigfopen(&[out], 1)); // 1: one output signal
    setsampfreq(250.0);
    for k in 0..500 {
        putvec(&[k]);
    }
    println!("{}", newheader("ramp")); // 0: ramp.hea is written
    let ai = Anninfo { name: b"qrs".to_vec(), stat: WFDB_WRITE };
    annopen("ramp", &[ai], 1);
    for t in [100, 350] {
        let a = Annotation { time: t, anntyp: NORMAL as i8, ..Default::default() };
        putann(0, &a);
    }
    wfdbquit(); // completes ramp.dat and ramp.qrs
}
\end{lstlisting}

% ==========================================================================================
\section{The command-line tools in Rust}
\label{sec:rust-tools}

The 33 command-line tools (\code{rdsamp}, \code{rdann}, \code{wrsamp}, \code{xform},
\ldots) are separate programs of the crate, one per file of \file{rust/src/bin}, built in
\file{rust/target/release/} (or installed in \file{\textasciitilde/.cargo/bin} by
\cmd{cargo install}), and run by their names. Each follows its C program line by line.
The tools, their options, their output and their exit status are those of the C tools,
and are documented in \cref{chap:tools}.

\begin{lstlisting}[style=shell]
export PATH=$PWD/rust/target/release:$PATH
export WFDBX_INPUT_DIR=tests/data     # run from the root of the repository
rdsamp -r 100s -t 10 -p     # the first 10 seconds of 100s, in physical units
rdann -r 100s -a atr        # its annotations
rdsamp --legacy -r 100s     # legacy mode: --legacy as the first argument
\end{lstlisting}

There is no launcher: a tool is run by its own name, and it has no \opt{-{}-version}
option, like the C tools. The tools have the same names as the original WFDB tools: if
both are installed, the one that comes first in the \code{PATH} runs. When an output file
or the standard output cannot be written (a full disk, for example), a tool reports it
(``PROGRAM: write error on standard output'') and exits with status 1; with
\opt{-{}-legacy}, it keeps the original exit status. A tool whose standard output is a
closed pipe ends with status 141, as a C program ended by \code{SIGPIPE}.

\paragraph{Environment variables.} The tools, like the rest of the library, use the
environment variables of \cref{tab:rust-environment}, also for the files that they read
or create themselves (\cref{sec:tool-files}): the \opt{-i} file of \cmd{wrsamp},
\cmd{edf2mit}, \cmd{wav2mit}, \cmd{a2m} and \cmd{ad2m}, the input files of
\cmd{ahaecg2mit} and the \opt{-S} script of \cmd{xform} are read from the input directory,
and the WAV file of \cmd{mit2wav}, the EDF file of \cmd{mit2edf}, the master header of
\cmd{snip -m} and the \file{.des}, \file{.ecg}, \file{.key} and \file{.txt} files of
\cmd{parsescp} are created in the output directory. Such a file cannot be read from a URL.
In legacy mode, the files that a tool reads itself are opened as named, and the output
directory, if set, is a plain prefix of every output name.

\begin{table}[htbp]
\centering
\caption{The environment variables used by the Rust library}
\label{tab:rust-environment}
\begin{tabularx}{\linewidth}{@{}lL@{}}
\toprule
\tabhead Variable & \tabhead Meaning \\
\midrule
\code{WFDBX\_INPUT\_DIR} & the input directory (a directory, a \code{\textasciitilde}
  path or a URL), when \code{wfdbx\_set\_input\_dir} does not set one \\
\code{WFDBX\_OUTPUT\_DIR} & the output directory, when \code{wfdbx\_set\_output\_dir} does
  not set one \\
\code{WFDBX\_NETWORK} & network access is on if it is a nonzero number, when
  \code{wfdbx\_set\_network} does not set it \\
\code{WFDBANNSORT} & \code{0} to keep annotation files in the order written (no sorting
  when they are closed) \\
\code{WFDBGVMODE} & the default getvec mode (\code{setgvmode} with a negative mode) \\
\code{WFDBCAL} & the calibration file read by \code{calopen(NULL)} (default
  \file{wfdbcal}) \\
\code{WFDB} & the WFDB path, used in legacy mode only (outside legacy mode, it is
  ignored, with a message) \\
\code{HOME} (\code{USERPROFILE} on Windows) & the home directory, for
  \code{\textasciitilde} paths \\
\bottomrule
\end{tabularx}
\end{table}

Tools that write signal parameters require them as options: \code{wrsamp} (with
\opt{-o}) requires \opt{-F} (frequency), \opt{-G} (gain), \opt{-R} (ADC resolution) and,
for calibrated signals, \opt{-U} (units); \code{wav2mit} requires \opt{-G}, and \opt{-U}
unless the gain is 0; \code{ad2m} and \code{ahaecg2mit} require \opt{-F}, \opt{-G},
\opt{-R}, and \opt{-U} unless the gain is 0; \code{a2m}, \code{rdann}, \code{ann2rr},
\code{sortann} and \code{m2a} require \opt{-F} when the record's header gives no
frequency (\cref{chap:tools}).

% ==========================================================================================
\section{Testing}
\label{sec:rust-testing}

The tests of the Rust implementation are the unit tests of the modules, the integration
tests of \file{rust/tests} and the examples of the doc comments (which read the sample
record \code{100s} from \file{tests/data}). From the directory \file{rust}:

\begin{lstlisting}[style=shell]
cargo test --release                       # all the tests, and the doc examples
cargo test --release --test metadata       # one file of rust/tests
cargo test --release shared_cases          # selected tests, by name
cargo clippy --release --all-targets       # lints
\end{lstlisting}

The integration tests cover the WFDB library functions (\file{signals.rs},
\file{annotations.rs}, \file{wfdbio.rs}, \file{bulk.rs}, the C standard functions in
\file{cstd.rs}), finding files and the settings (\file{files.rs}), network access
(\file{network.rs}), the high-level functions, the metadata functions (\file{metadata.rs},
with the cases shared by the seven languages, in \file{tests/metadata}) and
\code{check\_record} (\file{check.rs}). They run in temporary copies of the test records
of \file{tests/data}. The tests of network access use a local HTTP server on 127.0.0.1
only, and need \cmd{curl}. The tests of \code{check\_record} write the corpus of
\file{python/tests/check\_corpus.py} with \cmd{python3} (or the program named by the
environment variable \code{PYTHON}), NumPy and the Python package of \file{python/src},
and compare the reports with its \file{expected.txt}.
\cmd{cargo run -{}-release -{}-example check\_cases DIR} prints the reports for a corpus in
\file{DIR}.

The cross-language conformance suite runs every tool and compares its results, byte for
byte, with those of the original C programs (1135 cases), and, with \opt{-{}-legacy}, in
legacy mode (883 cases). It runs the tools of \file{rust/target/release}:

\begin{lstlisting}[style=shell]
cargo build --release --manifest-path rust/Cargo.toml
PYTHONPATH=$PWD/python/src python3 tests/conformance/harness.py check --impl rust -j8
PYTHONPATH=$PWD/python/src python3 tests/conformance/harness.py check --impl rust -j8 --legacy
\end{lstlisting}
